% ------------------------------------------------------------------
%  Registro de decisiones de arquitectura
%  Fragmento del Subdocumento 4.2, traido desde contenido.tex con
%  \parte{partes/14_m_decisiones_adr}. Las rutas son relativas a la carpeta del
%  subdocumento, nunca a la de main.tex.
% ------------------------------------------------------------------

\chapter{Registro de decisiones de arquitectura}
\label{sec:m-registro-de-decisiones-de-arquitectura}

Registro de las quince decisiones de arquitectura en formato MADR, con contexto cuantificado, alternativas evaluadas, criterio de selección, decisión y consecuencias. Es entregable contractual y se mantiene actualizado durante toda la ejecución (RT-02.04).

El registro es entregable contractual y se mantiene actualizado durante toda la ejecución (RT-02.04). Cubre las quince decisiones que condicionan la arquitectura de esta propuesta: estilo arquitectónico, conectividad, modelo de despliegue, persistencia, mensajería, identidad, movilidad, destino del sistema de almacenes heredado, recuperación ante desastres, almacenamiento local, intercambio con el canal moderno, absorción del peak estacional, puerta de enlace de servicios, observabilidad y gestión de secretos.

\paragraph{Finalidad.}
Cada decisión se fundamenta con el mismo rigor que un informe de ingeniería: contexto cuantificado desde el caso, alternativas reales de mercado evaluadas, criterio de selección técnico-económico ---costo total de propiedad, equipo de TI de cuatro personas del CLIENTE, latencia y resiliencia---, decisión precisa y sus contrapartidas con mitigaciones. La trazabilidad usa los códigos exactos de las Bases Técnicas Transversales (RT-cc.nn), de las Bases Administrativas (Art. n\textdegree{}) y de los capítulos del Caso 02. Cuando una decisión de arquitectura resuelve una de las decisiones del registro del subdocumento 3, se indica su número; las dieciséis primeras de ese registro corresponden al numeral 16.1 del caso y se citan como tales.

\section{Índice de decisiones}
\label{sub:indice-de-decisiones}

\begin{tablalafrox}{Índice del registro de decisiones de arquitectura}{tab:t104}%
  {G{0.10} Y G{0.11} F{0.26}}%
  {\cab{ADR} & \cab{Título} & \cab{Estado} & \cab{Decisión relacionada}}
  ADR-01 & Estilo arquitectónico: monolito modular (Django 5.x LTS / Python 3.12) & Aprobado & D5 (AL v1) \\
  ADR-02 & Conectividad y redundancia WAN: Starlink LEO + SD-WAN & Aprobado & Decisión 16.1 N\textdegree{} 26 \\
  ADR-03 & Modelo de despliegue híbrido (Art. 16) & Aprobado & D7 (AL v1) \textperiodcentered{} Decisión 16.1 N\textdegree{} 17 \\
  ADR-04 & Persistencia políglota por dominio (CAP) & Aprobado & D2/D4 (AL v1) \textperiodcentered{} Decisión 16.1 N\textdegree{} 18 \\
  ADR-05 & Mensajería asíncrona: RabbitMQ local + SQS FIFO/EventBridge & Aprobado & D4 (AL v1) \\
  ADR-06 & Identidad híbrida (Modelo B): Keycloak & Aprobado & D6 (AL v1) \textperiodcentered{} Decisión 16.1 N\textdegree{} 34 \\
  ADR-07 & Movilidad de terreno: nativa Android (Kotlin) & Aprobado & Decisión de proyecto N\textdegree{} 19 \\
  ADR-08 & Destino del WMS legado de 2013 & Aprobado & Decisión 16.1 N\textdegree{} 14 \\
  ADR-09 & DR: activo-pasivo warm standby multi-región & Aprobado & Decisión 16.1 N\textdegree{} 27 \\
  ADR-10 & Almacenamiento on-premise y niveles RAID & Aprobado & Decisión 16.1 N\textdegree{} 28 \\
  ADR-11 & Integración B2B/EDI: hub GS1 con capa anticorrupción & Aprobado & RF-12 \textperiodcentered{} RF-01.10 \\
  ADR-12 & Absorción del peak de septiembre (cómputo elástico) & Aprobado & Decisión 16.1 N\textdegree{} 30 \\
  ADR-13 & Puerta de enlace de servicios: Amazon API Gateway & Aprobado & D8 (AL v6.2) \\
  ADR-14 & Observabilidad: plataforma única con emisión on-premise & Aprobado & D14 (AL v6.2) \\
  ADR-15 & Gestión de secretos: servicio administrado & Aprobado & D13 (AL v6.2) \textperiodcentered{} SEC-01 \\
\end{tablalafrox}

\section{ADR-01 --- Estilo Arquitectónico: Monolito Modular vs. Microservicios}
\label{sub:adr-01-estilo-arquitectonico-monolito-modu}

Estado: Aprobado

Fecha: 2026-09-05

Decisiones relacionadas: D5 (Arquitectura Lógica, 2026-09-03, confirmada 09-04); D8 (Amazon API Gateway, actualizada 2026-09-05 --- ver ADR-13); D14 (observabilidad de plataforma única, 2026-09-06 --- ver ADR-14).

Trazabilidad normativa: BTT RT-02.02 (arquitectura modular, fronteras explícitas, despliegue independiente de componentes críticos) \textperiodcentered{} BTT sección 2.3 (la sofisticación no reemplaza la pertinencia: microservicios sin volumen que los justifique = error de ingeniería) \textperiodcentered{} RT-02.01 (8 capas) \textperiodcentered{} RT-02.03 (ISO/IEC/IEEE 42010) \textperiodcentered{} RT-09.05 (cuello de botella) \textperiodcentered{} BA Art. 16 (híbrido) \textperiodcentered{} Caso 02 Cap. 14.2/15 (volumetría) \textperiodcentered{} RNF-19.01--04 (capacidad, escalado, cuello de botella, pruebas 1,5$\times$).

\paragraph{Contexto y problema.}
La carga del caso es transaccional y concentrada, pero no plana: \textasciitilde{}31.000 pedidos/mes y 260.000 líneas/mes sobre 14.200 clientes y 8.400 productos, con un peak de \textasciitilde{}105 TPS en la ventana de despacho (dimensionamiento: \textasciitilde{}20 TPS confirmación de picking con 120 terminales, \textasciitilde{}8 TPS preventa, \textasciitilde{}35 TPS movimientos + sincronización, burst $\times$3). Pérdidas operacionales traducibles a dinero: conteo cíclico con 2,3 \% de diferencia, merma por vencimiento 1,7 \% del inventario, OTIF 82,4 \% contra meta 95 \% y \$4,2 millones/mes de descuadre de caja. El equipo de TI del CLIENTE tiene 4 personas para operar la solución completa durante 36 meses de Operación.

\paragraph{Alternativas evaluadas.}
\begin{tablalafrox}{ADR-01 \textperiodcentered{} Alternativas evaluadas}{tab:t105}%
  {F{0.14} Y F{0.14} F{0.14} F{0.14} F{0.14}}%
  {\cab{Alternativa} & \cab{Descripción} & \cab{Soporte a escala (\textasciitilde{}105 TPS)} & \cab{Equipo de 4 TI} & \cab{Complejidad de operación} & \cab{Rechazo/Aceptación}}
  A. Monolito modular (Django 5.x LTS, Python 3.12) & Único despliegue de proceso con módulos internos (apps Django por dominio: WMS, preventa, distribución, trazabilidad, portal), Amazon API Gateway en la frontera (ADR-13), workers Celery para asincronía & Suficiente: 6 tareas Fargate en peak absorben 105 TPS con headroom $\times$1,5 & 1 framework, 1 lenguaje, 1 artefacto CI/CD; curva de reposición corta & Baja: un despliegue por ambiente, rollback simple, observabilidad OTel única & Ganadora \\
  B. Microservicios en Kubernetes (EKS) & 20+ servicios acoplados a lo operacional (bodega, reparto, preventa, telemetría, facturación), mesh Istio, 3+ nodos por AZ & Sobra capacidad operativa para 105 TPS & Inviable: cluster operado por 4 personas (9--13 planos de control, versionado de contratos, backing services por dominio) & Alta: 2-3 SRE dedicados, curva de despliegue semanal, lock-in de operadores & Rechazada (BTT sección 2.3: volumen no la justifica) \\
  C. Monolito "clásico" sin límites de módulos (refactor del WMS 2013) & Continuidad sobre el monolito actual sin fronteras de contexto ni despliegue independiente & Cumple hoy, no escala a multi-sitio & N/D (proveedor desaparecido, sin soporte) & Alta (deuda estructural) & Rechazada (viola RT-02.02) \\
\end{tablalafrox}

\paragraph{Criterio de selección y justificación técnico-económica.}
Escala real (BTT sección 2.3): para 31.000 pedidos/mes y 105 TPS de diseño, el costo de coordinación de 20 servicios supera cualquier ganancia. El umbral donde los microservicios se justifican (típicamente $>$10\textsuperscript{3}--10\textsuperscript{4} TPS, equipos de 5+ por dominio) está dos órdenes de magnitud por encima del caso.

Equipo de 4 personas (RT-02.02/RT-02.03): un monolito modular tiene 1 artefacto, 1 pipeline, 1 proceso: la reposición de personal y la mantención son lineales. EKS exige 9--13 componentes de plano de control (Kube-apiserver, etcd, Ingress, autoscaler, RBAC, CNI, mesh\dots{}) más 20 servicios; el costo fijo de operarlo supera al de la nube usada.

Despliegue independiente (RT-02.02): la modularidad del monolito se materializa en procesos y tareas separables ---tarea portal (canal moderno, hito enero 2029), workers Celery (celery-reconciliation, celery-erp-sync), shipper de colas--- que escalan y se despliegan de forma independiente sin particionar el dominio.

TCO (56 meses): A y B convergen en cómputo AWS, pero B agrega \textasciitilde{}3 nodos EKS dedicados ($\approx$ USD 1.200--1.600/mes) + SRE; A escala con Fargate (2$\rightarrow$6 tareas) solo en septiembre (ADR-12). El diferencial B$-$A se estima en USD 40.000--70.000 en 56 meses por infraestructura y soporte, sin beneficio funcional neto.

Latencia/resiliencia: el monolito modular en nube + borde local (ADR-03) cumple RNF-05.01 ($\leq$1 s) y RNF-02.01 (24 h) porque lo que corre local es el mismo kernel WMS; particionar en servicios no aporta resiliencia adicional y añade saltos de red.

\paragraph{Decisión adoptada.}
Monolito modular en Django 5.x LTS sobre Python 3.12, con apps Django por dominio y límites de contexto explícitos (fundamentos de DDD y módulos de la Arquitectura Lógica (Subdocumento 4.1)), desplegado en ECS Fargate (tarea web + tarea portal + workers Celery), con Amazon API Gateway (D8 \textperiodcentered{} ADR-13), observabilidad OpenTelemetry sobre plataforma única en nube (D14 \textperiodcentered{} ADR-14) y borde local por sitio (ADR-03). Los módulos críticos (WMS offline, shipper de reconciliación, workers erp-sync, tarea portal) se despliegan de forma independiente en procesos separados, satisfaciendo RT-02.02 sin adoptar microservicios.

\paragraph{Consecuencias, contrapartidas y mitigaciones.}
\begin{tablalafrox}{ADR-01 \textperiodcentered{} Consecuencias, contrapartidas y mitigaciones}{tab:t106}%
  {F{0.26} F{0.26} Y}%
  {\cab{Consecuencia} & \cab{Trade-off} & \cab{Mitigación}}
  Riesgo de acoplamiento entre apps si no se respetan los límites & Deuda estructural intrínseca al monolito & Architecture Fitness Functions (ArchUnit/Django: contratos de módulo) en CI; revisión por módulo (RT-02.01) \\
  Un fallo de proceso puede amplificarse a todo el despliegue & Golpe de blast radius mayor que en microservicios & Separación en tareas/workers; health checks ALB; retry/backoff; observabilidad por transaction\_id (D9) \\
  Escalado vertical limitado del monolito & El peak se absorbe con más réplicas de la misma imagen & Auto-scaling predictivo+reactivo del peak de septiembre (ADR-12, 2$\rightarrow$6 tareas) \\
  Menor "trendiness" frente a la evaluación técnica & La BTT penaliza los microservicios injustificados según su sección 2.3 & Documentación explícita del análisis de escala/TCO (este ADR) como evidencia de pertinencia \\
\end{tablalafrox}

\section{ADR-02 --- Conectividad y Redundancia WAN: Starlink LEO + SD-WAN}
\label{sub:adr-02-conectividad-y-redundancia-wan-star}

Estado: Aprobado

Fecha: 2026-09-05

Decisión relacionada: Decisión N\textdegree{} 26 del registro de decisiones (Subdocumento 3) (redundancia de enlace de cada sitio).

Trazabilidad normativa: BTT RT-03.17 (enlace redundante con caminos y proveedores distintos, conmutación $\leq$ 5 min) \textperiodcentered{} RT-03.10 (autonomía 24 h) \textperiodcentered{} RT-03.19 (edge) \textperiodcentered{} RNF-13.01 (autonomía 24 h CD / 14 h terreno) \textperiodcentered{} RNF-13.07 (enlace redundante $\leq$ 5 min) \textperiodcentered{} RNF-13.08 (ancho de banda dimensionado) \textperiodcentered{} RT-10.05 (ventana crítica 05:30--07:00, indisponibilidad cero) \textperiodcentered{} Caso 02 (numerales 6.12 y 8: cortes de fibra, Concepción sin respaldo, Los Ángeles sin señal en madrugada) \textperiodcentered{} BA Art. 16.

\paragraph{Contexto y problema.}
La operación depende del enlace de datos en horarios en que no existe plan manual (Restricción N\textdegree{} 1: indisponibilidad cero 05:30--07:00). Verificado en el caso:

\begin{itemize}
  \item \textbf{Talca.} fibra se corta 4 veces al año sin enlace de respaldo; la bodega debe aguantar 24 h sin enlace.
  \item \textbf{Concepción.} no tiene respaldo alguno.
  \item \textbf{Cross-docks (Curicó, Chillán, Los Ángeles).} conectividad exclusivamente por red móvil; Los Ángeles pierde señal intermitentemente entre las 03:00--05:00, exactamente su ventana de operación de 3 h.
  \item \textbf{Impacto económico.} un día sin despacho de la mañana = operación caída (96 camiones).
\end{itemize}

\paragraph{Alternativas evaluadas.}
\begin{tablalafrox}{ADR-02 \textperiodcentered{} Alternativas evaluadas}{tab:t107}%
  {F{0.12} Y F{0.12} F{0.12} Y F{0.12}}%
  {\cab{Alternativa} & \cab{Descripción} & \cab{Cobertura Los Ángeles 03:00--05:00} & \cab{Cumple RT-03.17 ($\leq$5 min, caminos distintos)} & \cab{Costo 56 meses (referencial)} & \cab{Rechazo/Aceptación}}
  A. Fibra + LTE + Starlink LEO con SD-WAN (tri-camino) & Starlink como principal en cross-docks y respaldo en CDs; LTE dual (2 proveedores); SD-WAN con conmutación automática y BGP & Cobertura confirmada por constelación LEO (independiente de torres 4G) & Sí: 3 caminos, 2 proveedores LTE distintos + satelital & CAPEX \$2.012.500 CLP + OPEX \$50.400.000 CLP (56 m) = \$52.412.500 CLP (\textasciitilde{}USD 58.236) (cotización 5 sitios) & Ganadora \\
  B. Fibra + LTE 4G puro & Redundancia solo terrestre (fibra + 4G/Cat-12 femto backhaul) & No: en Los Ángeles la torre 4G falla justo en la ventana & Parcial: dos caminos pero el radio comparte tecnología física de torres & Sin CAPEX satelital; LTE puro \textasciitilde{}USD 20--30 K en 56 m & Rechazada (no cierra el caso Los Ángeles ni el "sin red de respaldo" de Concepción) \\
  C. Satélite GEO tradicional (VSAT) & Enlace de órbita geoestacionaria & Latencia 500--700 ms de ida y vuelta; débil en lluvia & Sí en camino, pero RTT penaliza RT-03.11/RNF-05.01 & Superior a LEO (equipos terminales GEO) & Rechazada (latencia y costo) \\
\end{tablalafrox}

\paragraph{Criterio de selección y justificación técnico-económica.}
Riesgo funcional: el único escenario que B no resuelve ---señal móvil intermitente en Los Ángeles 03:00--05:00--- es el que el caso declara crítico. La alternativa A aporta un camino físicamente independiente de las torres 4G (constelación LEO) y de las trincheras de fibra.

RT-03.17/RNF-13.07: conmutación automática $\leq$ 5 min por diseño SD-WAN (detección de pérdida $<$ 5 s, failover $<$ 30 s según práctica de diseño del sitio); la redundancia de caminos es de física distinta (fibra $\rightarrow$ satelital $\rightarrow$ LTE), no solo de proveedor.

Ventana crítica (RT-10.05): el despacho 05:30--07:00 no depende de un único medio; el tri-camino convive con la autonomía local de 24 h (RNF-13.01) para el caso de pérdida total (RT-03.10).

TCO: la oferta Starlink (5 sitios, 56 meses) es USD \textasciitilde{}58.236; frente a LTE puro que no cumple el requisito, el costo se justifica como prima de resiliencia del servicio crítico (costo de un día sin despacho >> costo del enlace). FinOps: CAPEX unitario bajo (\$350.000 CLP/kit), reposición 15 \% por ciclo de vida (RT-08.13).

Equipo de 4 TI: Starlink Enterprise gestionado (cero mantención de infraestructura de radio), SD-WAN con políticas centralizadas; sin ingeniería de redes satelitales in-house.

\paragraph{Decisión adoptada.}
Tri-camino WAN por sitio con SD-WAN: (1) fibra (enlace principal en Talca y Concepción), (2) Starlink Enterprise LEO ---principal en los 3 cross-docks (Curicó, Chillán, Los Ángeles) y respaldo automático en los 2 CD--- y (3) LTE dual (2 proveedores) con módulo 4G Cat-12. Conmutación automática $<$ 30 s entre caminos (sección de arquitectura de despliegue, red privada virtual IPsec) y con planes/precios declarados en la oferta (T-11 C27 \textperiodcentered{} T-12 RT-08.13). Los cross-docks salen por Starlink directo a SQS/IoT/SSM (configuración crítica) y sincronizan detalle a Talca por AMQPS entre brokers (RT-03.11).

\paragraph{Consecuencias, contrapartidas y mitigaciones.}
\begin{tablalafrox}{ADR-02 \textperiodcentered{} Consecuencias, contrapartidas y mitigaciones}{tab:t108}%
  {F{0.26} F{0.26} Y}%
  {\cab{Consecuencia} & \cab{Trade-off} & \cab{Mitigación}}
  Costo recurrente de planes Starlink & La prima de conectividad satelital es el mayor gasto de red & Planes tier por sitio (1 TB CDs / 500 GB cross-docks); revisión trimestral FinOps (RT-03.06) \\
  Skyline/interferencia puntual en cielos despejados (lluvia fuerte) & Degradación momentánea de throughput & SD-WAN con conmutación por política; LTE como tercer camino independiente \\
  Dependencia de tarifas/dólar en 56 meses (TCM USD 900 CLP) & Variación cambiaria en OPEX & Cotización con supuesto de reposición y cláusula de revisión de tarifas; registro en supuestos \\
  Sostenimiento de 5 kits (RT-08.13) & Gestión de repuestos & 15 \% de reposición cotizado; ciclo de vida declarado en la sección de Emplazamiento e inventario de este documento \\
\end{tablalafrox}

\section{ADR-03 --- Modelo de Despliegue Híbrido: Borde Operacional On-Premise vs. Nube AWS}
\label{sub:adr-03-modelo-de-despliegue-hibrido-borde-}

Estado: Aprobado

Fecha: 2026-09-05

Decisiones relacionadas: D7 (Subdocumento 4.1); Decisión N\textdegree{} 17 del registro de decisiones (Subdocumento 3) (asignación de componentes).

Trazabilidad normativa: BA Art. 16.1--16.4 (híbrido obligatorio, carga principal en nube, criterios 16.2, exigencias 16.3/16.4) \textperiodcentered{} RT-03.01 (región primaria/secundaria) \textperiodcentered{} RT-03.02 (multi-AZ) \textperiodcentered{} RT-03.10 (autonomía) \textperiodcentered{} RT-03.19 (edge) \textperiodcentered{} RNF-02.01 (24 h) \textperiodcentered{} RNF-05.01 ($\leq$ 1 s en cámara) \textperiodcentered{} RNF-05.02 (guantes $-$22 \textdegree{}C) \textperiodcentered{} Caso 02 Cap. 6 (cámara sin cobertura) y Cap. 16.1.

\paragraph{Contexto y problema.}
Tres hechos confluyen: (a) el Art. 16 exige carga principal en nube pública + componentes on-premise (inadmisible solo-nube o solo-on-prem); (b) la cámara de congelado a $-$22 \textdegree{}C no tiene cobertura de señal y el picking exige $\leq$ 1 s de confirmación (RNF-05.01), inviable con RTT de 40--80 ms hacia la nube más procesamiento remoto; (c) la bodega debe operar 24 h sin enlace (RNF-02.01/RT-03.10). El perfil de carga es nocturno (22:00--06:00) y de madrugada (05:30--07:00).

\paragraph{Alternativas evaluadas.}
\begin{tablalafrox}{ADR-03 \textperiodcentered{} Alternativas evaluadas}{tab:t109}%
  {Y Y F{0.12} F{0.12} F{0.12} F{0.12}}%
  {\cab{Alternativa} & \cab{Descripción} & \cab{RT-03.10 (24 h)} & \cab{RNF-05.01 ($\leq$1 s)} & \cab{Art. 16.1 (carga en nube)} & \cab{Rechazo/Aceptación}}
  A. Híbrido: WMS maestro on-prem + borde por sitio + nube AWS (carga principal) & WMS y BD transaccional en Talca (maestro), borde edge en Concepción y cross-docks, nube para OLTP cloud, canal moderno, analítica, IoT, respaldo inmutable & Sí (autonomía local 24 h CD / 14 h terreno) & Sí (confirmación en el dispositivo local) & Sí & Ganadora \\
  B. Solo nube (WMS en AWS) & Toda la operación en la nube & No: sin enlace, bodega muere en minutos & No: RTT + procesamiento no sostiene $\leq$1 s en cámara & No (no es híbrida) & Inadmisible (BA Art. 16) \\
  C. Solo on-premise (WMS maestro + todo local) & Centro de datos propio con todo el cómputo & Sí & Sí & No (no es híbrida) & Inadmisible (BA Art. 16) \\
\end{tablalafrox}

\paragraph{Criterio de selección y justificación técnico-económica.}
Cumplimiento normativo: solo A satisface Art. 16 en sus cuatro numerales (16.1 híbrido, 16.2 justificación por componente, 16.3 exigencias de nube, 16.4 exigencias on-premise). La tabla maestra de emplazamiento es la sección (a) de este documento (36 componentes).

Latencia (RNF-05.01): el borde local ejecuta la confirmación de picking en el propio sitio; la alternativa nube añade $\geq$40--80 ms RTT y procesamiento remoto, por encima del umbral de 1 s en la cámara.

Autonomía (RNF-02.01/RT-03.10): el diseño del caso exige supervivencia sin enlace; el maestro local con réplica cloud (DMS CDC, RPO $\leq$ 15 min) da continuidad (ADR-09) sin depender de la WAN durante el corte.

TCO/equipo: el cómputo local se acota a lo imprescindible (23 componentes on-premise/híbrido), el resto se delega a servicios administrados AWS (ADR-12) --- menor costo fijo que un DC propio y menor riesgo que una operación 100 \% nube.

Carga principal en nube (Art. 16.1): canal moderno (portales), OLTP cloud, ingesta IoT, analítica/BI, respaldo inmutable y DRP viven en AWS (sa-east-1, DRP us-east-1 --- RT-03.01).

\paragraph{Decisión adoptada.}
Arquitectura híbrida obligatoria con borde operacional on-premise: WMS maestro + BD transaccional en Talca (VM-01/VM-02), edge autónomo en Concepción (VM-C01/VM-C02), mini-WMS en cross-docks (E-01), y máxima utilización de AWS administrado (ECS Fargate, Aurora, DynamoDB, IoT Core, S3, Redshift/QuickSight) para la carga principal. Detalle por componente en la sección de modelo de emplazamiento híbrido.

\paragraph{Consecuencias, contrapartidas y mitigaciones.}
\begin{tablalafrox}{ADR-03 \textperiodcentered{} Consecuencias, contrapartidas y mitigaciones}{tab:t110}%
  {Y F{0.26} Y}%
  {\cab{Consecuencia} & \cab{Trade-off} & \cab{Mitigación}}
  Dos dominios operativos que mantener (on-prem + nube) & Mayor superficie que la nube pura & IaC total (RT-03.03), SSM Agent (gestión sin puertos entrantes), Zero Trust (D7) \\
  El maestro local exige administración de BD/hardware & Costo de operación on-premise & Dimensionamiento con headroom $\times$1,5 y RAID 10 (ADR-10); NOC 24$\times$7 (RNF-21.01) \\
  La identidad debe sobrevivir offline & Despliegue de caché local de solo lectura & Modelo B (ADR-06), sin maestro local a promover \\
  La réplica cloud aumenta costos de AWS & DRP por réplica continua & Aurora replicada solo para lo crítico (A-01/A-02); respaldos 3-2-1-0 (ADR-09) \\
\end{tablalafrox}

\section{ADR-04 --- Persistencia Políglota por Dominio}
\label{sub:adr-04-persistencia-poliglota-por-dominio}

Estado: Aprobado

Fecha: 2026-09-05

Decisiones relacionadas: D2 (series de tiempo condicionadas al muestreo, consolidado en OLAP); D4 (nube AWS + PostgreSQL/Redis/Aurora/DynamoDB/S3+Redshift); Decisión N\textdegree{} 18 del registro de decisiones (Subdocumento 3) (motor por dominio, RT-05.02).

Trazabilidad normativa: BTT RT-05.02 (paradigma, garantías transaccionales y posición CAP por dominio) \textperiodcentered{} RT-05.01 (diccionario de datos) \textperiodcentered{} RT-05.11--05.15 (migración/volúmenes) \textperiodcentered{} RNF-09.01 (retiro $<$ 2 h) \textperiodcentered{} RNF-09.02 (retención 5 años) \textperiodcentered{} RNF-11.02 (OLAP desacoplado del OLTP) \textperiodcentered{} RSA D.S. 977/96 (retención legal de trazabilidad) \textperiodcentered{} Caso 02 Cap. 16.1/16.2 y Cap. 4.9.

\paragraph{Contexto y problema.}
Cada dominio tiene exigencias distintas de consistencia, volumen y retención: pedidos/inventario/trazabilidad demandan ACID y reconciliación determinista tras cortes; la ingesta de sensores y telemetría es escritura masiva de baja latencia; las series de tiempo (temperatura de cámara y camión) requieren compresión y consulta por rango; la analítica exige columnar separado para no degradar picking; y la retención sanitaria impone 5 años de trazabilidad (piso legal) con 6 meses adicionales por vida útil. Un único motor no puede cumplir las cuatro posiciones CAP sin sacrificar desempeño.

\paragraph{Alternativas evaluadas.}
\begin{tablalafrox}{ADR-04 \textperiodcentered{} Alternativas evaluadas}{tab:t111}%
  {F{0.19} F{0.19} F{0.16} F{0.19} Y}%
  {\cab{Dominio} & \cab{Alternativa A} & \cab{Alternativa B} & \cab{Alternativa C} & \cab{Ganador}}
  Transaccional (pedidos, inventario, trazabilidad, facturación) & PostgreSQL ACID (CP) & MySQL & MongoDB (AP) & A --- PostgreSQL (CP: consistencia inmediata exigida por operación offline y reconciliación) \\
  OLTP cloud + réplica DRP WMS & Aurora PostgreSQL & RDS MySQL & DocumentDB & Aurora (administrado, failover $<$30 s, PITR 35 días, réplica WAL) \\
  IoT raw (senores, telemetría) & DynamoDB (AP) & PostgreSQL puro & Kafka+Elasticsearch & DynamoDB (escritura serverless $<$10 ms p99, TTL 30 días) \\
  Series de tiempo (frío) & Capa OLAP: S3 Parquet + Redshift Serverless (columna) & InfluxDB & Particionado manual & OLAP (S3+Redshift) (adoptado según muestreo --- D2; sin segundo motor; consolidado 5 años) \\
  Analítica / BI / lake 5 años & S3 + Redshift Serverless + QuickSight & Snowflake & Athena sobre S3 puro & S3+Redshift+QuickSight (columnar desacoplado, RNF-11.02) \\
  Retención legal / respaldo inmutable & S3 Object Lock / Glacier + Backup Vault & Cinta/tape local & NAS local puro & S3 Object Lock/Glacier (inmutabilidad WORM, custodia RT-06.26+27) \\
\end{tablalafrox}

\paragraph{Criterio de selección y justificación técnico-económica.}
RT-05.02 y CAP: la posición se declara por dominio: transaccional CP (consistencia inmediata; la operación offline y la reconciliación vuelta a línea lo exigen), IoT raw AP (eventual, TTL corto), series de tiempo AP columnar en la capa OLAP, analítica AP columnar. Declarar la posición evita el error de elegir un motor "para todo".

Volumen (RT-05.15): \textasciitilde{}31.000 pedidos/mes, 260.000 líneas, 2,4 M unidades y telemetría de 28 sensores + 18 termógrafos $\rightarrow$ persistencia manejable en PostgreSQL (índices + PostGIS para georreferencia de rutas) con DynamoDB absorbiendo el pico de IoT sin saturar el transaccional (RNF-11.02).

Trazabilidad (RNF-09.01/09.02, D.S. 977/96): el registro de eventos por lote (EPCIS GS1, Decisión N\textdegree{} 35 del registro de decisiones (Subdocumento 3)) se conserva en PostgreSQL (5 años) y su respaldo inmutable replica el piso legal; la réplica Aurora da RPO $\leq$15 min.

Equipo de 4 TI: tres motores administrados o de bajo mantenimiento (PostgreSQL/Aurora, DynamoDB, Redshift Serverless) --- sin motores dedicados exóticos (InfluxDB descartado: segundo motor a operar, D2). La serie de tiempo consolidada no requiere un motor adicional: vive en la capa OLAP (S3 Parquet $\rightarrow$ Redshift).

\paragraph{Decisión adoptada.}
Persistencia políglota por dominio: PostgreSQL+PostGIS (transaccional WMS on-premise, CP), Aurora PostgreSQL (OLTP cloud + réplica DRP, failover $<$30 s), DynamoDB (IoT raw, TTL 30 días, AP), S3 (lake analítico 5 años) con Redshift Serverless y QuickSight (OLAP, incluye la serie de tiempo consolidada de temperatura), y S3 Object Lock/Glacier + Backup Vault para retención legal inmutable (RSA D.S. 977/96). La matriz CAP está declarada y trazada a RT-05.02 en este documento (sección de Datos, ADR-04).

\paragraph{Consecuencias, contrapartidas y mitigaciones.}
\begin{tablalafrox}{ADR-04 \textperiodcentered{} Consecuencias, contrapartidas y mitigaciones}{tab:t112}%
  {Y F{0.26} Y}%
  {\cab{Consecuencia} & \cab{Trade-off} & \cab{Mitigación}}
  Consistencia fuerte en el transaccional limita escalado horizontal de escritura & El volumen (\textasciitilde{}105 TPS) no lo justifica & Réplicas de lectura (2 Aurora readers) y partición lógica por instalación (RF-02.01/02.02) \\
  DynamoDB eventual abre ventanas de lectura inconsistente cortas & Telemetría no es crítico de negocio & TTL 30 días + consolidado en la capa OLAP (S3 Parquet/Redshift, 5 años) con validación de rango \\
  Amazonas de datos (S3) requiere gobierno & Costo y gobernanza del lake & Glue + catálogo y QuickSight con modelos aprobados; FinOps (RT-03.06) \\
  Integración entre motores exige ETL & Latencia de analítica != OLTP & Batch nocturno y distribución de datos por dominio (RT-05.01) \\
\end{tablalafrox}

\section{ADR-05 --- Mensajería Asíncrona: Broker Local (RabbitMQ) + Colas Nube (SQS FIFO / EventBridge)}
\label{sub:adr-05-mensajeria-asincrona-broker-local-r}

Estado: Aprobado

Fecha: 2026-09-05

Decisión relacionada: D4 (Subdocumento 4.1); shipper en VM-03 (D-AL-03).

Trazabilidad normativa: BTT RT-02.06 (idempotencia) \textperiodcentered{} RT-02.07 (deduplicación) \textperiodcentered{} RT-03.12 (sincronización/reconciliación tras reconexión) \textperiodcentered{} RT-03.10 (autonomía) \textperiodcentered{} RNF-07.01 (sincro $<$ 10 min) \textperiodcentered{} RNF-13.01 (24 h) \textperiodcentered{} Zero Trust (BA Art. 21) \textperiodcentered{} Caso 02 Cap. 16.1 (reglas de reintento/devolución).

\paragraph{Contexto y problema.}
El caso exige orden estricto y sin pérdida en la cadena pedido$\rightarrow$despacho$\rightarrow$cobranza (una línea de preventa fuera de orden rompe la reconciliación determinista), pero la operación transcurre offline: bodega 24 h, terreno 14 h. Además, todo el tráfico debe ser outbound (Zero Trust: sin conexiones entrantes al on-premise). REST síncrono no sobrevive a un corte a mitad de ventana; Kafka aporta orden y durabilidad pero exige operación y no resuelve el broker local por sitio.

\paragraph{Alternativas evaluadas.}
\begin{tablalafrox}{ADR-05 \textperiodcentered{} Alternativas evaluadas}{tab:t113}%
  {F{0.12} F{0.12} F{0.12} Y Y F{0.12}}%
  {\cab{Alternativa} & \cab{Orden estricto} & \cab{Buffer offline 24 h} & \cab{Zero Trust (outbound)} & \cab{Equipo de 4 TI} & \cab{Rechazo/Aceptación}}
  A. RabbitMQ local + SQS FIFO/EventBridge (nube), shipper & SQS FIFO por group\_id (pedido/entidad) garantiza orden & RabbitMQ 24 h (\textasciitilde{}3 M mensajes) en cada sitio & Shippper VM-03 publica por HTTPS/443 (VPC Endpoint SQS), sin entrantes & RabbitMQ es maduro y liviano; SQS/EventBridge administrados & Ganadora \\
  B. Kafka (clúster local + MSK nube) & Orden por partición & Posible local, pero requiere clúster mínimo 3 brokers por sitio & Requiere conectividad zookeeper/plano de control & Elevada: 3+ brokers por sitio + MSK; operación de retención y rebalances & Rechazada (escala y operación no la justifican) \\
  C. REST síncrono puro & N/D & No: cae con la WAN & Parcial & Baja & Rechazada (viola RT-03.10/RNF-13.01) \\
\end{tablalafrox}

\paragraph{Criterio de selección y justificación técnico-económica.}
Resiliencia (RT-03.10/RNF-13.01): el broker local retiene transacciones y telemetría durante el corte y las reproduce al reconectar; el shipper publica en orden cronológico e idempotente (idempotency keys, RT-02.06/02.07) a la cola FIFO.

Reconciliación determinista (RT-03.12): celery-reconciliation procesa la cola al recuperar el enlace; claves de idempotencia evitan duplicados (RT-02.07) y el orden de publicación por entidad garantiza consistencia en el maestro.

Zero Trust (BA Art. 21): SQS FIFO se consume via VPC Endpoint SQS/PrivateLink (HTTPS/443); AMQPS 5671 queda solo entre brokers on-premise (cross-dock $\rightarrow$ Talca). No se abren puertos entrantes.

TCO/equipo: RabbitMQ (este) + SQS/EventBridge (administrados) $>$ Kafka (+MSK) para 105 TPS; el costo operativo de operar clústeres Kafka en 5 sitios es desproporcionado para 4 personas.

\paragraph{Decisión adoptada.}
Capa de mensajería híbrida: RabbitMQ (A-03) como broker de colas offline en cada sitio (buffers 24 h), con shipper en VM-03 (modo wms\_only) que publica el buffer hacia SQS FIFO (VPC Endpoint / HTTPS 443) con idempotencia; EventBridge como bus de eventos de negocio (eventos que disparan ERP-sync, alertas, auditoría) y AWS IoT Core (MQTT) para la ingesta edge (ADR-04/ADR-12). Workers celery-reconciliation y celery-erp-sync procesan en nube. Declarado en la sección de arquitectura de integración (subsección de mensajería) y en el ADR-05.

\paragraph{Consecuencias, contrapartidas y mitigaciones.}
\begin{tablalafrox}{ADR-05 \textperiodcentered{} Consecuencias, contrapartidas y mitigaciones}{tab:t114}%
  {F{0.26} F{0.26} Y}%
  {\cab{Consecuencia} & \cab{Trade-off} & \cab{Mitigación}}
  Dos brokers (local/nube) con estrategias distintas & Más superficie que un solo bus & Contratos de eventos versionados (AsyncAPI, RT-05.16/17); matriz de integración (vista de integración de este documento) \\
  Posible duplicado en el flujo local$\rightarrow$nube & Garantías "at-least-once" & Idempotency keys + procesamiento deduplicado (RT-02.06/07); buffer con ACK solo tras commit exitoso \\
  Volume de mensajes en septiembre (2$\times$) & Crecimiento de cola & SQS y broker dimensionados; auto-escalado de workers (ADR-12) \\
  Kafka descartado & Menor capacidad de reprocesamiento histórico & S3 como foso de eventos (reproceso por replays controlados, ADR-04) \\
\end{tablalafrox}

\section{ADR-06 --- Identidad y Autenticación Híbrida (Modelo B): Keycloak IdP Maestro en Nube + Caché Local de Solo Lectura}
\label{sub:adr-06-identidad-y-autenticacion-hibrida-m}

Estado: Aprobado

Fecha: 2026-09-05

Decisiones relacionadas: D6 (Subdocumento 4.1); Decisión N\textdegree{} 34 del registro de decisiones (Subdocumento 3) (credenciales de usuarios externos sin correo).

Trazabilidad normativa: BTT RT-12.10 (aprovisionamiento $\leq$ 24 h) \textperiodcentered{} RT-12.11 (autenticación en perfil operacional: guantes, rotación 38 \%, dispositivos compartidos) \textperiodcentered{} RT-12.12 (credenciales de externos sin correo) \textperiodcentered{} RT-03.10/RNF-13.01 (offline 24 h CD / 14 h terreno) \textperiodcentered{} BA Art. 22 (MFA para acceso externo) \textperiodcentered{} RF-15.01--15.05 (identidad centralizada) \textperiodcentered{} RF-06.08 (OTP conductores externos $\leq$ 2 min) \textperiodcentered{} Decisión N\textdegree{} 34 del registro de decisiones (Subdocumento 3).

\paragraph{Contexto y problema.}
La operación exige autenticación incluso sin enlace: la bodega trabaja 24 h y el terreno 14 h desconectados, con 62 preventistas y \textasciitilde{}200 conductores (de 10 transportistas externos) sin correo corporativo, dispositivos compartidos entre turnos y rotación de personal de preparación de 38 \% anual. Un IdP solo-nube paraliza la bodega ante un corte; un maestro local tradicional (AD) duplica la administración de identidad y crea un maestro a promover.

\paragraph{Alternativas evaluadas.}
\begin{tablalafrox}{ADR-06 \textperiodcentered{} Alternativas evaluadas}{tab:t115}%
  {Y F{0.17} F{0.17} F{0.17} Y}%
  {\cab{Alternativa} & \cab{Offline 24 h/14 h} & \cab{Externos sin correo (OTP)} & \cab{Un solo maestro (sin promoción)} & \cab{Rechazo/Aceptación}}
  A. Keycloak IdP maestro en nube (ECS Fargate) + caché local de solo lectura TTL 8 h + offline tokens (8 h/14 h) & Sí (validación local de firma OIDC) & Sí (OTP emitido por jefe de flota, RF-06.08) & Sí (Modelo B, D6) & Ganadora \\
  B. IdP nube puro (ej. Cognito-only) & No: sin red, sin autenticación & Parcial (depende de conectividad) & Sí & Rechazada (paraliza CD/terreno, RT-03.10) \\
  C. Active Directory tradicional on-premise & Sí & Parcial (sin flujo OTP para externos sin dominio) & No: maestro local a promover en DR & Rechazada (complejidad de federación + DR de identidad) \\
  D. Cognito como dependencia & No & Parcial & --- & Rechazada (D6: Cognito solo alternativa; descendencia de dependencia) \\
\end{tablalafrox}

\paragraph{Criterio de selección y justificación técnico-económica.}
Autonomía (RT-03.10/RNF-13.01): la caché local (A-05: VM-05 Talca, VM-C03 Concepción) valida localmente la firma OIDC y mantiene sesiones hasta 24 h CD / 14 h terreno; el token offline se renueva al inicio del turno en cobertura, de modo que la ventana offline siempre cubre la jornada (8 h bodega / 14 h reparto-preventa) sin autenticación en ruta.

Autoridad única (D6): todas las escrituras viven en el IdP maestro de la nube; los cambios de alta/baja/roles se propagan por export/import del Realm cifrado vía S3 (outbound) con $\Delta$ $\leq$ 8 h (TTL) y baja efectiva SCIM $\leq$ 24 h (RT-12.10). No existe maestro local a promover $\Rightarrow$ DR de identidad = misma autoridad en nube (no depende de switchover).

Externos (RT-12.11/12.12): OTP de un solo uso para conductores (RF-06.08, $\leq$ 2 min, sin correo) y credencial inicial en sesión de dotación; MFA para todo acceso externo y privilegiado (Art. 22).

Equipo de 4 TI: Keycloak administrado (Fargate) + federación OIDC con Amazon API Gateway (ADR-13); sin sincronización de directorio propietario.

\paragraph{Decisión adoptada.}
Modelo B de identidad (D6): Keycloak IdP maestro en AWS (ECS Fargate, 2 tareas Multi-AZ, backend Aurora) con caché local de solo lectura TTL 8 h en VM-05/VM-C03 y offline access tokens por perfil (8 h bodega, 14 h reparto/preventa). Sincronización del Realm cifrado por S3 (outbound, Zero Trust); MFA OIDC para externos; OTP para conductores externos. Declarado en la sección de arquitectura de seguridad (subsección de identidad, acceso y sesiones, ADR-06).

\paragraph{Consecuencias, contrapartidas y mitigaciones.}
\begin{tablalafrox}{ADR-06 \textperiodcentered{} Consecuencias, contrapartidas y mitigaciones}{tab:t116}%
  {F{0.26} F{0.26} Y}%
  {\cab{Consecuencia} & \cab{Trade-off} & \cab{Mitigación}}
  Revocación no es inmediata en el borde ($\Delta$ $\leq$ 8 h TTL) & Ventana de revocación por TTL & SCIM $\leq$ 24 h; baja reforzada por matrícula de activos y MDM (bloqueo de terminal) \\
  El Realm viaja por S3 (outbound) & Exposición del cifrado durante el traslado & Cifrado KMS del objeto; solo lectura de cachés; sin claves privadas locales \\
  OTP para externos depende del jefe de flota & Fricción en cambio sin aviso & Emisión previa a la ventana; listado de contingencia; rotación diaria (RF-12.21) \\
  Falla del IdP maestro deja sin altas/bajas & Degradación administrativa & Identidad en nube multi-AZ + us-east-1 DRP (ADR-09); las evaluaciones locales siguen operando \\
\end{tablalafrox}

\section{ADR-07 --- Estrategia de Movilidad de Terreno: Aplicación Nativa Android (Kotlin)}
\label{sub:adr-07-estrategia-de-movilidad-de-terreno-}

Estado: Aprobado

Fecha: 2026-09-05

Decisión relacionada: Decisión de proyecto N\textdegree{} 19 (nativa vs. híbrida vs. PWA).

Trazabilidad normativa: BTT RT-12.11 (autenticación en perfil operacional) \textperiodcentered{} RT-13.08 (interfaz de terreno: guantes térmicos, $-$22 \textdegree{}C, lluvia, 34 \textdegree{}C, una mano, turno sin conexión) \textperiodcentered{} RT-03.19 (edge) \textperiodcentered{} RNF-05.02 (guantes $-$22 \textdegree{}C) \textperiodcentered{} RNF-05.03 (curva de aprendizaje $\leq$ 2 h) \textperiodcentered{} RNF-06.01 (14 h sin señal) \textperiodcentered{} RF-03.02 (stock offline) \textperiodcentered{} RF-06.11/06.12/06.13 (QR/foto/impresora térmica) \textperiodcentered{} Caso 02 Cap. 6 (cámara sin cobertura, turnos de 30 min).

\paragraph{Contexto y problema.}
El terreno es el corazón de la operación: 62 preventistas, \textasciitilde{}200 conductores y 120 terminales de bodega operando con guantes térmicos a $-$22 \textdegree{}C, a una mano durante la descarga, bajo lluvia, sol directo y 14 h sin señal. Requiere escáner GS1, impresora térmica Bluetooth y POS móvil (PAX), y curva de aprendizaje $\leq$ 2 h. Una PWA no controla el SDK de escáner Zebra de forma confiable ni persiste un turno completo sin capa nativa; una híbrida (Flutter) agrega un runtime intermedio que degrada la interacción con periféricos industriales.

\paragraph{Alternativas evaluadas.}
\begin{tablalafrox}{ADR-07 \textperiodcentered{} Alternativas evaluadas}{tab:t117}%
  {F{0.12} Y Y F{0.12} F{0.12} F{0.12}}%
  {\cab{Alternativa} & \cab{Periféricos Zebra (DataWedge)/BT} & \cab{Offline 14 h} & \cab{UX guantes/1 mano/$-$22 \textdegree{}C} & \cab{Curva $\leq$ 2 h} & \cab{Rechazo/Aceptación}}
  A. Nativa Android (Kotlin), SQLite/Room + SDK Zebra DataWedge + scanner API & Total (intent com.symbol.datawedge) & SQLite con sync deduplicada (RF-03.17) & Full control de la UI (botones grandes, táctil frente a guantes) & SÍ (flujos guiados) & Ganadora \\
  B. PWA (canal web) & Limitado (Web Bluetooth parcial, sin DataWedge nativo) & Solo Service Worker con cuotas y poca fiabilidad para cámara/escáner & Media & Parcial & Rechazada para terreno; adoptada para autoatención del cliente (portal) \\
  C. Híbrida Flutter/React Native & Puentes parciales, mantención de capa nativa & Posible pero con runtime intermedio & Media & Media & Rechazada (complejidad y rendimiento de periféricos) \\
\end{tablalafrox}

\paragraph{Criterio de selección y justificación técnico-económica.}
Periféricos (RT-13.08/RF-06.13): DataWedge (intents) entrega disparo de escáner GS1, CW (camera wedge) para QR (RF-06.11) y salidas a impresora BW Zebra --- integraciones nativas estables que un runtime híbrido/PWA no garantiza.

Offline total (RNF-06.01/RT-03.19): SQLite/Room persiste el turno completo (pedidos, POD con foto/firma, cobros, devoluciones) y sincroniza con deduplicación e idempotencia al reconectar ($<$ 10 min, RNF-07.01); STock indicativo con antigüedad visible (RF-02.05) y resolución cronológica de conflictos (RF-02.05e).

Condiciones extremas (RNF-05.02/RT-12.11): la UI nativa permite contraste alto, botones grandes, entrada por guantes (toque grueso), sin gestos complejos; dispositivos Rugged Zebra (EC55, TC58e, MC9400 Cold Storage) con Android 13+ y baterías de cámara/5G.

TCO: una misma app nativa cubre preventa, reparto y bodega (parque de terreno de la sección de modelo de emplazamiento híbrido); el costo de mantención de la plataforma nativa se asume como trade-off frente a las exigencias RNF (documentado en este ADR-07).

\paragraph{Decisión adoptada.}
Aplicación nativa Android (Kotlin) para preventa, reparto/cobranza y picking/recepción de bodega, con persistencia local SQLite/Room, SDK Zebra DataWedge (escáner GS1 + CW QR) y módulo de impresión Bluetooth (ZQ620) y POS (PAX). PWA (web) solo para la autoatención del canal moderno (portal, ADR-11). Dispositivos Rugged: EC55 (preventa), TC58e (reparto), MC9400 Cold Storage (cámara $-$22 \textdegree{}C), DS2208 (andén), penalizados por guantes/una mano (RNF-05.02).

\paragraph{Consecuencias, contrapartidas y mitigaciones.}
\begin{tablalafrox}{ADR-07 \textperiodcentered{} Consecuencias, contrapartidas y mitigaciones}{tab:t118}%
  {Y F{0.26} Y}%
  {\cab{Consecuencia} & \cab{Trade-off} & \cab{Mitigación}}
  Dos stacks de UI (nativa Android + portal web) & Mayor mantención que una sola tecnología & Automatización de CI/tests (Appium), contratos API compartidos; la web cubre el canal sin offline \\
  Versionado de app en parque de \textasciitilde{}270 dispositivos (62 EC55 preventa + \textasciitilde{}200 TC58e reparto + reservas; 174 MC9400 bodega) & Coordinación de despliegues de terreno & MDM (Gestión de Endpoints, F-03) con controles de distribución; actualizaciones en ventana sin operación \\
  SQLite frente a sincronización concurrente & Conflictos de stock & Claves de idempotencia + reconciliación cronológica (ADR-05) \\
  Dependencia de kits Android fabricantes (Zebra) & Ciclo de vida OS & Dispositivos con Android 13$\rightarrow$18 Garage/Zebra; política de intercambio RT-08.13 \\
\end{tablalafrox}

\section{ADR-08 --- Destino del WMS Legado de 2013}
\label{sub:adr-08-destino-del-wms-legado-de-2013}

Estado: Aprobado

Fecha: 2026-09-05

Decisión relacionada: Decisión N\textdegree{} 14 del numeral 16.1 del caso (reemplazar el WMS de 2013).

Trazabilidad normativa: RT-05.11--RT-05.15 (plan de migración por dominio, volúmenes y reversión) \textperiodcentered{} RT-03.10/RNF-02.01 (24 h) \textperiodcentered{} RNF-05.01 ($\leq$1 s) \textperiodcentered{} Caso 02, numeral 16.1, decisión N\textdegree{} 14 (el CLIENTE delega expresamente la decisión) y Cap. 9 (dolores: conteo 2,3 \%, merma 1,7 \%, rasgos de WMS 2013) \textperiodcentered{} estrategia azul-verde y reversión.

\paragraph{Contexto y problema.}
El WMS de 2013 corre en Talca con proveedor desaparecido (sin soporte, sin roadmap), planillas impresas (conteo cíclico con 2,3 \% de diferencia y ajustes sin investigación de causa), y no soporta multi-sitio ni operación de 24 h sin enlace. El CLIENTE delega expresamente en el proponente la decisión (numeral 16.1 del caso, decisión N\textdegree{} 14). Extenderlo arrastra el riesgo de fecha conocida: sin soporte, un incidente en la ventana crítica (05:30--07:00) no tendría plan B.

\paragraph{Alternativas evaluadas.}
\begin{tablalafrox}{ADR-08 \textperiodcentered{} Alternativas evaluadas}{tab:t119}%
  {F{0.13} Y G{0.10} Y F{0.13} F{0.13}}%
  {\cab{Alternativa} & \cab{Multi-sitio (RF-02.01/02.02)} & \cab{Offline 24 h} & \cab{Soporte y riesgo} & \cab{Costo} & \cab{Rechazo/Aceptación}}
  A. Reemplazo total en la Etapa 1 & Sí (topología configurable, consolidación multi-nodo) & Sí (RNF-02.01) & Proveedor nuevo con soporte contractual; conversión controlada & Mayor CAPEX, menor riesgo de operación & Ganadora \\
  B. Mantener/extender el WMS 2013 & No (arquitectura monocliente sin multi-sitio) & No & Proveedor desaparecido: sin parches ni soporte, un fallo = cero plan B en ventana crítica & Bajo CAPEX, riesgo operacional alto e ilimitado & Rechazada \\
  C. Coexistencia temporal (WMS 2013 + nuevo en un sitio) & Parcial (solo mientras dure el puente) & --- & Migración en olas reduce el riesgo & Medio & Complemento de la estrategia de despliegue (azul-verde), no fin último \\
\end{tablalafrox}

\paragraph{Criterio de selección y justificación técnico-económica.}
Requisitos (RF-02.x): slotting ABC, conteo cíclico ciego, picking FEFO dirigido, SSCC GS1, trazabilidad por lote --- el WMS 2013 no los contempla ni los puede incorporar sin proveedor.

Riesgo operacional: un sistema sin soporte expuesto a la ventana de indisponibilidad cero (RT-10.05) es un riesgo con fecha de impacto indefinido; el análisis de riesgos del caso lo califica como inaceptable frente al CAPEX del reemplazo.

Migración (RT-05.11--15): por dominio y en ventanas sin operación: maestros completos, ventas/pedidos 3 años, inventario 2 años, trazabilidad 5 años, CxC abierta; convivencia con el ERP/GDE (que no se reemplaza) vía capa anticorrupción (ADR-11).

Reversión y despliegue: estrategia azul-verde por ola y sitio, con umbrales de marcha blanca por oleada (RT-20.02/20.04) y plan de reversión para volver al WMS 2013 sin pérdida en caso de no superar los indicadores.

TCO 56 meses: el costo de reemplazo se paga en la Etapa 1 y se amortiza con la reducción de 2,3 \%$\rightarrow$\textasciitilde{}0,3 \% de diferencia de conteo y 1,7 \%$\rightarrow$<1 \% de merma (objetivos de RF/planilla del Cap. 18) y con OTIF 82,4 \%$\rightarrow$95 \%.

\paragraph{Decisión adoptada.}
Reemplazo total del WMS 2013 en la Etapa 1 por el módulo WMS del monolito modular (ADR-01) desplegado en Talca (maestro), Concepción (edge) y cross-docks (mini-WMS E-01), con migración por dominio (RT-05.11--15), oleadas por sitio y plan de reversión azul-verde documentado. El ERP administrativo/GDE no se reemplaza y se integra por la capa anticorrupción (ADR-11).

\paragraph{Consecuencias, contrapartidas y mitigaciones.}
\begin{tablalafrox}{ADR-08 \textperiodcentered{} Consecuencias, contrapartidas y mitigaciones}{tab:t120}%
  {F{0.26} F{0.26} Y}%
  {\cab{Consecuencia} & \cab{Trade-off} & \cab{Mitigación}}
  Conversión de datos en ventanas sin operación & Ventana de corte por dominio & Migración nocturna por dominio con verificación de conciliación (RT-05.12/05.14) \\
  Curva de aprendizaje del personal de bodega & Riesgo operacional en cutover & Marcha blanca por ola con certificación (RT-20.02/20.04); curva $\leq$ 2 h (RNF-05.03) \\
  El ERP sigue existiendo & Dualidad de maestros durante la cohabitación & Capa anticorrupción (ADR-11) y catálogo de equivalencias (RF-12.03) \\
  Costo de migración y reversiones & CAPEX de transición & Plan de reversión por ola; pruebas de restauración mensuales (RNF-20.07) \\
\end{tablalafrox}

\section{ADR-09 --- Estrategia de Recuperación ante Desastres (DR): Warm Standby Activo-Pasivo Multi-Región}
\label{sub:adr-09-estrategia-de-recuperacion-ante-des}

Estado: Aprobado

Fecha: 2026-09-05

Decisión relacionada: Decisión N\textdegree{} 27 del registro de decisiones (Subdocumento 3) (modalidad activo-pasivo, RTO/RPO, cadencia de pruebas).

Trazabilidad normativa: RT-07.01 (declarar y justificar modalidad activo-activo vs. activo-pasivo) \textperiodcentered{} RT-07.07 (pruebas reales $\geq$ 2 veces/año con informe y medición de RTO/RPO) \textperiodcentered{} RNF-20.06 (RTO $\leq$ 4 h, RPO $\leq$ 15 min) \textperiodcentered{} RNF-20.07 (3-2-1-0) \textperiodcentered{} BA Art. 20 (continuidad y pruebas semestrales) \textperiodcentered{} DRP local Talca$\rightarrow$Concepción (RTO +1--2 h) \textperiodcentered{} Decisión N\textdegree{} 27 del registro de decisiones (Subdocumento 3).

\paragraph{Contexto y problema.}
La ventana de despacho (05:30--07:00) no admite plan manual; una falla mayor del sitio primario debe recuperar el servicio con RTO $\leq$ 4 h y RPO $\leq$ 15 min (RNF-20.06) y probarse con conmutación real al menos 2 veces al año (RT-07.07/Art. 20). El sitio primario es Talca (WMS maestro on-premise); el respaldo secundario se apoya en la nube (replicación AWS) y en el borde autónomo de Concepción como DRP local.

\paragraph{Alternativas evaluadas.}
\begin{tablalafrox}{ADR-09 \textperiodcentered{} Alternativas evaluadas}{tab:t121}%
  {F{0.12} Y F{0.12} F{0.12} F{0.12} Y}%
  {\cab{Alternativa} & \cab{RTO/RPO} & \cab{Complejidad operativa} & \cab{Costo} & \cab{Cumple RT-07.01 (justificación)} & \cab{Rechazo/Aceptación}}
  A. Activo-pasivo warm standby multi-región (Aurora réplica DRP + DMS CDC) & RTO $\leq$ 4 h / RPO $\leq$ 15 min & Media: promote documentado, pruebas semestrales & Réplica + DRP regional (us-east-1) & Sí & Ganadora \\
  B. Activo-activo (multi-maestro) & RPO \textasciitilde{}0 & Alta: conflictos de escritura y reconciliación continua entre nodos & 2$\times$ infraestructura transaccional + reconciliación & Justificable pero no necesario al volumen & Rechazada (costo/complejidad desproporcionados, RT-07.01) \\
  C. Backup en frío (bajo demanda) & RTO $>$ 24 h (restauración completa) & Baja & Baja & No cumple RTO $\leq$ 4 h & Rechazada \\
  D. Continuidad intrarregional en sa-east-1 (sin us-east-1) & Falla de AZ o corrupción: RTO $\leq$ 4 h / RPO $\leq$ 15 min (multi-AZ + CDC) \textperiodcentered{} caída de toda la región sa-east-1: RTO 24--72 h desde respaldo inmutable (S3 Object Lock/Backup Vault) & Media: restauración y rehidratación; cada 6--12 meses requiere prueba real & Réplica y respaldo solo dentro de sa-east-1; menor costo de salida que A & No cumple RTO $\leq$ 4 h ante falla regional (RNF-20.06) & Contingencia condicionada: solo si el CLIENTE rechaza la transferencia internacional (Art. 23) --- no sustituye a A como plan base \\
\end{tablalafrox}

\paragraph{Criterio de selección y justificación técnico-económica.}
RTO/RPO (RNF-20.06): la réplica continua (DMS CDC) y la réplica de Aurora (WAL) mantienen RPO $\leq$ 15 min; el promote a la región DRP (us-east-1) más el DRP local Talca$\rightarrow$Concepción (procedimiento documentado de 5 pasos, RTO +1--2 h) cumplen los cuatros horas.

RT-07.01: activo-activo duplicaría la infraestructura transaccional y exigiría reconciliación de doble escritura (\textasciitilde{}105 TPS peak) sin beneficio medible frente al RTO comprometido; la modalidad activo-pasiva es proporcional al volumen (CAP consistente, ADR-04).

Pruebas (RT-07.07/Art. 20): conmutación real semestral con informe de RTO/RPO efectivos y plan de corrección de brechas; consistente con la retención sanitaria y la operación continua.

Identidad (ADR-06): el DR de identidad no exige conmutación local (el maestro está en la nube y las cachés A-05 siguen validando firmas): se elimina una clase entera de riesgos de DR.

Respaldo 3-2-1-0 (RNF-20.07): copia primaria on-premise + réplica local (Concepción) + pierna cloud (S3 Object Lock/Backup Vault, inmutable) + offsite + prueba mensual de restauración y retención legal (ADR-04).

Falla de AZ vs. caída de región (alternativa D): si el CLIENTE no aprueba la transferencia a us-east-1 (Art. 23, sección de arquitectura de seguridad, subsección de datos personales, residencia y transferencia internacional), la postura mínima es la continuidad intrarregional en sa-east-1: uso de la tercera AZ (sa-east-1a/1b/1c) y respaldo inmutable regional, que mantiene RTO $\leq$ 4 h / RPO $\leq$ 15 min ante falla de una zona de disponibilidad o corrupción de datos --- lo que el multi-AZ ya brinda ---, pero no ante una caída de toda la región sa-east-1, escenario en que la reconstrucción desde el respaldo inmutable toma 24--72 h e incumple RNF-20.06. Esta alternativa no usa los sitios on-premise del CLIENTE para la carga cloud (Talca/Concepción alojan solo el dominio on-premise, con su DRP local RTO +1--2 h); y es precisamente esta degradación la que motiva solicitar la aprobación de us-east-1 con resguardos, conforme al Art. 23.

\paragraph{Decisión adoptada.}
DR activo-pasivo warm standby multi-región: replicación continua del WMS (VM-02) hacia Aurora (ACM us-east-1 global replica) con DMS CDC, RPO $\leq$ 15 min y RTO $\leq$ 4 h; DRP local Talca$\rightarrow$Concepción (VM-C01/VM-C02) con RTO +1--2 h para la bodega; pruebas reales 2$\times$/año (RT-07.07) y respaldo 3-2-1-0 admitido por S3 Object Lock. Documentado en la sección de arquitectura de despliegue (subsección de recuperación ante desastres, ADR-09). Si el CLIENTE no aprueba la transferencia internacional (Art. 23), rige la alternativa D del citado ADR-09: continuidad intrarregional en sa-east-1 con tercera AZ y respaldo inmutable, aceptando el RTO de 24--72 h ante caída de toda la región, que es el impacto evaluado en el análisis del ADR-09.

\paragraph{Consecuencias, contrapartidas y mitigaciones.}
\begin{tablalafrox}{ADR-09 \textperiodcentered{} Consecuencias, contrapartidas y mitigaciones}{tab:t122}%
  {F{0.26} F{0.26} Y}%
  {\cab{Consecuencia} & \cab{Trade-off} & \cab{Mitigación}}
  Capacidad de cómputo ociosa en la región DRP & Costo de la réplica caliente & Aurora multi-AZ compartida con OLTP cloud; DRP us-east-1 solo componentes críticos \\
  Ventana de pérdida $\leq$ 15 min & No cero-datos & RPO $\leq$ 15 min aceptado (RNF-20.06); reconciliación de vuelta (RT-03.12) \\
  Pruebas de conmutación real & Riesgo de fallo durante la prueba & Ventanas de prueba con plan de reversión; indicadores medidos y corrección de brechas (RT-07.07) \\
  Replicación CDC sobre enlaces WAN & Costo de ancho de banda & WAL/DMS comprimido; ventana de replicación nocturna junto a sync batch \\
\end{tablalafrox}

\section{ADR-10 --- Almacenamiento On-Premise y Niveles RAID}
\label{sub:adr-10-almacenamiento-on-premise-y-niveles}

Estado: Aprobado

Fecha: 2026-09-05

Decisión relacionada: Decisión N\textdegree{} 28 del registro de decisiones (Subdocumento 3) (declarar y justificar nivel RAID frente a alternativas).

Trazabilidad normativa: BTT RT-03.14 (tolerancia a falla de disco y nivel RAID declarado) \textperiodcentered{} RNF-13.03 (equipos críticos redundantes) \textperiodcentered{} RNF-13.04 (tolerancia a falla de al menos 1 disco) \textperiodcentered{} RNF-13.05 (justificación del nivel RAID) \textperiodcentered{} Decisión N\textdegree{} 28 del registro de decisiones (Subdocumento 3) y Cap. 10 del caso (ventana crítica sin falla).

\paragraph{Contexto y problema.}
La BD transaccional del WMS (VM-02) concentra \textasciitilde{}105 TPS de diseño con escritura fuerte de sincronización y trazabilidad; la evidencia fotográfica del POD, los logs y el contenido de cámara suman datos de escritura secuencial menos críticos. Una falla de disco en la ventana de despacho (05:30--07:00) no puede detener la operación. El nivel RAID debe tolerar el fallo de al menos un disco (RNF-13.04) y justificarse (RNF-13.05/RT-03.14).

\paragraph{Alternativas evaluadas.}
\begin{tablalafrox}{ADR-10 \textperiodcentered{} Alternativas evaluadas}{tab:t123}%
  {F{0.10} F{0.10} F{0.10} Y F{0.10} F{0.10} Y}%
  {\cab{Nivel} & \cab{Tolerancia} & \cab{Reescritura/desempeño} & \cab{Rebuild frente a URE} & \cab{Costo (utilizable)} & \cab{Aplicación} & \cab{Rechazo/Aceptación}}
  RAID 10 & N-1 discos por espejo ($\geq$ 1, hasta N/2) & Excelente para IOPS de escritura & Rápido y acotado & 50 \% (4--8 NVMe) & BD transaccional y réplica (VM-02, VM-C02, PMe VM-01) & Ganadora (transaccional) \\
  RAID 6 (doble paridad + hot-spare) & Hasta 2 discos simultáneos & Bueno para secuencial & Rebuild más largo, pero seguro ante segundo fallo & 2/N (p.ej. 2 de 8) & Evidencia POD (fotos/firmas), logs, rollups de cámara y mini-WMS de cross-dock & Ganadora (datos/evidencia) \\
  RAID 5 & 1 disco & Medio & Riesgo USB de URE en rebuild a gran capacidad + sin doble redundancia & 1/N (económico) & No aplica & Rechazada (no cumple RNF-13.04 robustamente en discos grandes y rebuild) \\
  JBOD/Sin RAID & 0 & N/D & N/D & --- & --- & Inadmisible (viola RNF-13.04) \\
\end{tablalafrox}

\paragraph{Criterio de selección y justificación técnico-económica.}
Transaccional (RNF-13.04, RT-03.14): RAID 10 sobre NVMe Superdome/Kioxia (sección de dimensionamiento y plan de capacidad, subsección de almacenamiento) entrega $>$100.000 IOPS frente a \textasciitilde{}840 IOPS necesarias (105 TPS $\times$ 8 E/S), con latencia de escritura determinista y reparación acotada --- crítico en la ventana de despacho.

Evidencia/logs: RAID 6 (doble paridad + hot-spare) para fotografías POD, logs y rollups; tolera el segundo fallo durante el rebuild, relevante en discos de alta capacidad y baja rotación de escritura.

Rechazo de RAID 5: en discos modernos de 8--16 TB, la probabilidad de URE (uncorrectable read error) durante un rebuild extenso no es despreciable; el caso exige tolerancia a fallo de al menos un disco en operación competitiva (RNF-13.04/05) y RAID 5 solo cubre uno.

Coherencia con Ceph (hipervisor): el pool Ceph N+1 de los nodos de cómputo complementa los niveles RAID de las VMs (redundancia de infraestructura, RNF-13.03); la pieza cross-dock usa contenedores con disco local del mini-PC recubierto por el respaldo nocturno (RNF-20.07).

TCO: el sobre-costo del 50 \% de RAID 10 solo se aplica a la BD crítica; el grueso del almacenamiento (evidencia, logs, telemetría) usa RAID 6 con solo 2 de 8 discos de paridad, manteniendo el costo de almacenamiento total contenidamente respecto de RAID 1/0 puro.

\paragraph{Decisión adoptada.}
RAID 10 para los servicios transaccionales del WMS (VM-02/VM-C02 y núcleo del maestro) con NVMe, y RAID 6 con hot-spare para evidencia fotográfica/POD, logs y rollups de cámara; RAID 5 descartado y justificado frente a alternativas (RNF-13.05/RT-03.14), con redundancia de equipos críticos por hipervisor Ceph N+1 (RNF-13.03). Registrado en la sección de dimensionamiento y plan de capacidad (subsección de almacenamiento) y en la sección de modelo de emplazamiento híbrido (ADR-10).

\paragraph{Consecuencias, contrapartidas y mitigaciones.}
\begin{tablalafrox}{ADR-10 \textperiodcentered{} Consecuencias, contrapartidas y mitigaciones}{tab:t124}%
  {F{0.26} F{0.26} Y}%
  {\cab{Consecuencia} & \cab{Trade-off} & \cab{Mitigación}}
  50 \% de pérdida útil en la BD & Mayor costo de almacenamiento crítico & Solo en la capa transaccional; resto en RAID 6 \\
  Rebuild de RAID 6 durante fallas & Ventana de vulnerabilidad al segundo fallo & Hot-spare + alertas NOC 24$\times$7 (RNF-21.01) y monitoreo SMART \\
  Complejidad de niveles mixtos & Operación de dos políticas & Política única de respaldo 3-2-1-0 (RNF-20.07) sobre ambas capas \\
  Dependencia del proveedor de hardware & Reposición de discos & SLA de reemplazo $\leq$ 4 h para componente crítico (sección de Emplazamiento de este documento); stock $\geq$ 10 \% parque \\
\end{tablalafrox}

\section{ADR-11 --- Integración B2B/EDI con Supermercados: Hub GS1 Centralizado con Capa Anticorrupción}
\label{sub:adr-11-integracion-b2b-edi-con-supermercad}

Estado: Aprobado

Fecha: 2026-09-05

Decisiones relacionadas: RF-12 (mensajería electrónica EDI); RF-01.10 (maestro interno de productos, tabla de equivalencias por cadena); A-04 (Capa Anticorrupción del ERP).

Trazabilidad normativa: BTT RT-05.23 (estándares sectoriales de intercambio --- EDI) \textperiodcentered{} RT-16.16 (documentos cifrados con integridad y retención) \textperiodcentered{} RT-16.17 (firma electrónica Ley N\textdegree{} 19.799) \textperiodcentered{} RNF-12.01 (canal moderno en producción antes de enero 2029) \textperiodcentered{} RF-12.03 (equivalencias GTIN por cadena) \textperiodcentered{} RF-12.06 (bandeja de excepciones) \textperiodcentered{} RF-01.10 --- estándares GS1 (EANCOM/GS1 XML, EPCIS) invocados por el Caso 02 Cap. 16.2 \textperiodcentered{} Caso 02 Cap. 9 (cadenas del canal moderno).

\paragraph{Contexto y problema.}
Los supermercados del canal moderno (p.ej. Walmart, Cencosud, SMU) exigen intercambio electrónico EDI dentro de sus ventanas de recepción (pedidos, envío de guías/facturas, acuse de recibo). Hoy el caso lo resuelve por correo/telefono y con diferencias de maestro por cadena (numeral 16.1 del caso, decisión N\textdegree{} 11). Conectores punto-a-punto por cadena multiplican los adaptadores, duplican la lógica de mapeo y disparan el mantenimiento ante cada cambio de especificación de la cadena; además, el ERP (que no se reemplaza) expone una frontera frágil.

\paragraph{Alternativas evaluadas.}
\begin{tablalafrox}{ADR-11 \textperiodcentered{} Alternativas evaluadas}{tab:t125}%
  {Y F{0.12} Y F{0.12} F{0.12} F{0.12}}%
  {\cab{Alternativa} & \cab{Estandarización GS1} & \cab{Esfuerzo por cadena nueva} & \cab{Aislamiento del ERP} & \cab{Madurez del ecosistema} & \cab{Rechazo/Aceptación}}
  A. Hub EDI centralizado (GS1) + Capa Anticorrupción (ACL) & Mapeo único a GS1 EANCOM/GS1 XML; EPCIS para trazabilidad & Conector configurado (perfiles por cadena), sin desarrollo a medida & El ERP solo habla con la ACL (nunca escritura directa) & Estándar de la industria; validación de esquemas & Ganadora \\
  B. Conectores punto-a-punto por cadena & Modelado ad-hoc por cadena & Desarrollo y pruebas por cada cadena & Frontera expuesta directamente al ERP & Baja & Rechazada (esfuerzo $\times$ cadenas y acoplamiento) \\
  C. Portal manual de documentos & No es EDI & Descarga/upload manual & N/D & Baja & Solo complemento (bandeja e impresos de respaldo) \\
\end{tablalafrox}

\paragraph{Criterio de selección y justificación técnico-económica.}
Estandarización (Cap. 16.2): el caso exige estándares GS1; el hub declara un modelo canónico (EANCOM D.01B/GS1 XML para pedidos/despachos/facturas, EPCIS para eventos de trazabilidad) y mapea cada cadena contra el modelo, no contra el ERP.

Esfuerzo marginal: una cadena nueva se agrega por configuración de perfil (equivalencias GTIN, RF-12.03) y pruebas de conexión, no por código; el hub centraliza la validación, el retry y la auditoría.

Aislamiento del ERP (ACL A-04): el ERP (que no se reemplaza, ADR-08) queda detrás de la capa anticorrupción: celery-erp-sync lee contratos OpenAPI del ERP y publica notificaciones por SQS; nunca hay escritura directa desde una cadena al ERP (Zero Trust, BA Art. 21).

Plazo contractual: el canal moderno (incluido EDI) entra en producción antes del hito enero 2029 (RNF-12.01), por lo que la Etapa 2 concentra portal y EDI detrás de la misma puerta de enlace (Amazon API Gateway, ADR-13) y del hub.

Equipo de 4 TI: el hub EDI se aloja como módulo del monolito (ADR-01) en la DMZ de AWS (N-01\dots{}N-03) con servicios administrados (API Gateway, SQS, EventBridge), sin adaptadores propietarios por cadena.

\paragraph{Decisión adoptada.}
Hub EDI centralizado basado en GS1 (EANCOM/GS1 XML para textos, EPCIS para trazabilidad) en la nube (DMZ, módulo de integraciones del monolito), con conector configurable por cadena, tabla de equivalencias GTIN/cadena (RF-12.03), bandeja de excepciones (RF-12.06) y Capa Anticorrupción (ACL, A-04) hacia el ERP/GDE. Canal moderno operativo $\leq$ enero 2029 (RNF-12.01). La evidencia POD se articula con la GDE y el acuse de recibo electrónico (RF-12.12/12.13, RT-16.17/16.18).

\paragraph{Consecuencias, contrapartidas y mitigaciones.}
\begin{tablalafrox}{ADR-11 \textperiodcentered{} Consecuencias, contrapartidas y mitigaciones}{tab:t126}%
  {F{0.26} F{0.26} Y}%
  {\cab{Consecuencia} & \cab{Trade-off} & \cab{Mitigación}}
  Dependencia de especificaciones por cadena & Cambios en el perfil de cada cadena & Versionado de contratos (AsyncAPI, RT-05.16/17) y pruebas de integración contínua contra los bancos de pruebas de cada cadena \\
  El ERP sigue siendo el emisor tributario & Latencia de la GDE/factura & Acuse dentro de los 30 min de la descarga (RF-12.13); la ACL abstrae el ERP \\
  Equivocaciones de GTIN por cadena & Diferencias de maestro & Equivalencias por cadena (RF-12.03) + bandeja de excepciones visible (RF-12.06) \\
  Costo de la infraestructura DMZ H2 & Recursos del hub & Servicios administrados y FinOps (RT-03.06); el hub comparte la tarea portal (ADR-01) \\
\end{tablalafrox}

\section{ADR-12 --- Absorción del Peak de Septiembre y Perfil de Carga No Plano}
\label{sub:adr-12-absorcion-del-peak-de-septiembre-y-}

Estado: Aprobado

Fecha: 2026-09-05

Decisión relacionada: Decisión N\textdegree{} 30 del registro de decisiones (Subdocumento 3) (cuello de botella del peak de septiembre y estrategia).

Trazabilidad normativa: BTT RT-09.05 (identificar el cuello de botella y cómo se detecta/resuelve) \textperiodcentered{} RT-03.06 (FinOps) \textperiodcentered{} RT-03.08 (instancias reservadas/ahorro) \textperiodcentered{} RT-03.09 (cómputo serverless para carga variable) \textperiodcentered{} RNF-19.01--04 (capacidad, escalado horizontal automático, cuello de botella, pruebas 1,5$\times$) \textperiodcentered{} Caso 02 Cap. 14.2/15 (perfil no plano; congelamiento 1--25 sept) \textperiodcentered{} RT-10.05 (congelamiento de cambios en septiembre/diciembre).

\paragraph{Contexto y problema.}
El perfil de carga no es plano: preventa 09:00--18:00, preparación 22:00--06:00, despacho 05:30--07:00 (96 camiones), sincronización de flota 17:00--20:00, y septiembre casi duplica el volumen durante 3 semanas (1.400 $\rightarrow$ 2.600 entregas/día). Sobredimensionar para el promedio diario es un error declarado por el propio caso; sobredimensionar estático para el peak paga 12 meses la capacidad de 3 semanas.

\paragraph{Alternativas evaluadas.}
\begin{tablalafrox}{ADR-12 \textperiodcentered{} Alternativas evaluadas}{tab:t127}%
  {Y F{0.17} Y F{0.17} F{0.17}}%
  {\cab{Alternativa} & \cab{Costo en régimen} & \cab{Comportamiento en peak} & \cab{Detección de cuello de botella} & \cab{Rechazo/Aceptación}}
  A. Cómputo elástico serverless/contenedores (Fargate 2$\rightarrow$6, Celery 2$\rightarrow$4, Lambda, Aurora escalada) & Pago por uso; escala solo cuando sube la carga & Auto-escalado predictivo (calendario sept/diciembre) + reactivo (CPU/cola) & Amazon CloudWatch / Prometheus; umbrales de cola y TPS visibles (RT-09.05) & Ganadora \\
  B. Sobredimensionamiento estático (12$\times$2 vCPU fijo + Aurora 2$\times$) & Pago fijo del 100 \% los 56 meses & Nunca degrada, pero se paga todo el año & Requiere monitoreo manual & Rechazada (costo: \textasciitilde{}40--60 \% más OPEX por el peak inactivo, contra FinOps RT-03.06) \\
  C. Overprovision moderado + colas más largas (sin escala automática) & Intermedio & Los workers cuelan y la reconciliación se desplaza a horario no crítico & Parcial & Complementa a A para cargas cola (reconciliación) \\
\end{tablalafrox}

\paragraph{Criterio de selección y justificación técnico-económica.}
Perfil no plano (Cap. 14.2/15): la capacidad se calcula al peak $\times$1,5 (RNF-19.04) = 3.900 entregas/día, con \textasciitilde{}105 TPS de diseño de despacho y pruebas de carga de 1,5$\times$peak $\approx$ 160 TPS / 5.850 entregas (RNF-19.04). Fargate 2$\rightarrow$6 tareas (Django), Celery 2$\rightarrow$4 y Lambda 100$\rightarrow$500 cubren la ratio de las 3 semanas sin pago residual.

Cuello de botella identificado (RT-09.05): confirmación/persistencia transaccional de pedidos e ingesta de telemetría de la flota. Se declara el punto de saturación y se vigila con métricas de negocio (OTIF, pedidos no preparados) y de infraestructura (cola SQS, CPU, TPS) --- RNF-19.03.

FinOps (RT-03.06) y RT-03.08/03.09: la base reservada (RI/Savings Plan) cubre la capacidad de régimen; el crecimiento del peak se paga con cómputo efímero (Fargate/Lambda). Ampliar Aurora a xlarge + 2 readers y DynamoDB con autoscaling absorbe el pico sin compra fija.

Congelamiento (RT-10.05): del 1 al 25 de septiembre y en diciembre no se despliegan cambios; la escala predictiva se programa por calendario (calendario de eventos) y se valida en la marcha blanca de septiembre del año 1.

Equipo de 4 TI: la operación es declarativa (IaC, RT-03.03): el mismo despliegue escala y decrece sin intervención manual; las políticas de alerta de desviación de presupuesto (RT-03.06) sustituyen el control manual del costo.

\paragraph{Decisión adoptada.}
Cómputo elástico con auto-escalado horizontal (Fargate Django 2$\rightarrow$6, Celery 2$\rightarrow$4, Lambda 100$\rightarrow$500, Aurora db.r6g.large$\rightarrow$xlarge +2 readers, DynamoDB on-demand) con escala predictiva por calendario (septiembre/diciembre) y reactiva (CPU, profundidad de cola, TPS); base de capacidad con Savings Plan/instancias reservadas de la carga de régimen (RT-03.08) y la diferencia del peak como cómputo efímero (RT-03.09). Monitoreo del cuello de botella declarado (RT-09.05/RNF-19.03) y pruebas 1,5$\times$peak (RNF-19.04).

\paragraph{Consecuencias, contrapartidas y mitigaciones.}
\begin{tablalafrox}{ADR-12 \textperiodcentered{} Consecuencias, contrapartidas y mitigaciones}{tab:t128}%
  {F{0.26} F{0.26} Y}%
  {\cab{Consecuencia} & \cab{Trade-off} & \cab{Mitigación}}
  Costo de proveer la parte efímera & Dependencia del autoscaler & 0 riesgo de pay-what-you-use: presupuestos y alarmas (RT-03.06); base RI para régimen \\
  La cola puede crecer en el burst & Latencia de reconciliación & Vúmenes dimensionados (SQS $\times$ 24 h buffer); workers autoescalados; congelamiento de cambios (RT-10.05) \\
  Complejidad del autoscaler predictivo & Configuración de calendario & Configuración IaC por evento (EventBridge); ensayo en marcha blanca de sept; aumentos conservadores $\times$1,5 \\
  Surgimiento de cuello de botella nuevo tras el peak & Identificación tardía & Monitoreo de síntomas de negocio (RF-16.03) + pruebas de estrés al punto de quiebre (RNF-19.04) \\
\end{tablalafrox}

\section{ADR-13 --- Puerta de Enlace de Servicios: Amazon API Gateway}
\label{sub:adr-13-puerta-de-enlace-de-servicios-amazo}

Estado: Aprobado

Fecha: 2026-09-05 (formalizado como ADR el 2026-09-06)

Decisión relacionada: D8 (Arquitectura Lógica (Subdocumento 4.1), Capa 3).

Trazabilidad normativa: BA Art. 21.2 (puerta de enlace con autenticación, autorización, cuotas, límites de tasa, validación de esquema e inspección de carga útil) \textperiodcentered{} RT-02.01 (capa 3 del modelo de referencia) \textperiodcentered{} RT-11.11 \textperiodcentered{} RT-05.16/05.18 (contratos y OAuth 2.1/mTLS) \textperiodcentered{} RT-02.02 (contratos versionados retro-compatibles) \textperiodcentered{} Art. 16.3 (preferencia por servicios administrados).

\paragraph{Contexto y problema.}
La Capa 3 es el punto por donde entra todo: las APIs de negocio, la sincronización diferida de preventa y reparto ---que es tráfico de primera clase, no un anexo--- y los tres portales de la DMZ. El Art. 21.2 no pide \guillemotleft{}un gateway\guillemotright{}: pide autenticación, autorización, cuotas, límites de tasa, validación de esquema e inspección de carga útil, todo en el borde. Y el CLIENTE opera con 4 personas de TI: cualquier componente que haya que parchar, dimensionar y sostener 24$\times$7 compite con la operación.

\paragraph{Alternativas evaluadas.}
\begin{tablalafrox}{ADR-13 \textperiodcentered{} Alternativas evaluadas}{tab:t129}%
  {F{0.12} F{0.12} F{0.12} Y Y F{0.12}}%
  {\cab{Alternativa} & \cab{Operación} & \cab{Integración con el borde} & \cab{Cuotas y esquema} & \cab{Costo para un equipo de 4} & \cab{Veredicto}}
  A. Amazon API Gateway (administrado) & Cero nodos que operar; escala con la carga & Nativa con CloudFront, WAF, Shield y ALB; autorizador OIDC contra Keycloak & Validación de esquema OpenAPI, cuotas y límites por cliente y por ruta, trazabilidad por transacción & Nulo en operación; costo por invocación declarado en FinOps & Ganadora \\
  B. Kong Gateway (open source, autoalojado) & Nodos propios con alta disponibilidad, parches y versiones a cargo del equipo & Requiere integrar manualmente con WAF/ALB & Equivalente vía plugins & Alto: un componente crítico más en la ruta de la venta, operado por 4 personas & Rechazada \\
  C. Sin puerta de enlace, exponiendo el ALB directo a los módulos & Mínima & --- & No cumple Art. 21.2 (sin cuotas ni validación de esquema en el borde) & --- & Rechazada (incumplimiento normativo) \\
\end{tablalafrox}

\paragraph{Criterio de selección y justificación técnico-económica.}
Cumplimiento literal del Art. 21.2 sin desarrollo propio: autenticación OIDC, autorización por rol, cuotas y límites de tasa por cliente, validación de esquema e inspección de carga útil son capacidades nativas.

Equipo de 4 personas (Cap. 2.4 del caso): la alternativa B agrega un componente crítico en la ruta de la venta que hay que dimensionar, parchar y recuperar. En la ventana de despacho 05:30--07:00, con indisponibilidad cero comprometida, ese es exactamente el riesgo que no conviene asumir.

Coherencia entre vistas: la puerta de enlace es la misma en la arquitectura lógica, en el modelo de emplazamiento híbrido, en la arquitectura de despliegue y en la estructura de costos de la oferta. El Art. 16.4 in fine califica de incoherencia grave que un componente declarado en una vista no aparezca en las demás, y por eso la decisión se toma una sola vez y se propaga.

Reversibilidad (Art. 16.3): los contratos son OpenAPI 3.1 y AsyncAPI 2.6, estándares abiertos; la lógica de negocio vive en Django, no en el gateway. Migrar a Kong u otro gateway no exige reescribir servicios, solo reconfigurar rutas y autorizadores. La dependencia es de configuración, no de código, y así queda declarada en la matriz de reversibilidad.

\paragraph{Decisión adoptada.}
Amazon API Gateway como Capa 3 única, con autorizador OIDC contra el Keycloak maestro, validación de esquema OpenAPI, cuotas y límites de tasa por actor y por ruta, versionado /v\{major\} y asignación de transaction\_id propagado a la Capa 8. Kong Gateway queda declarado como alternativa open source evaluada y no adoptada.

\paragraph{Consecuencias, contrapartidas y mitigaciones.}
\begin{tablalafrox}{ADR-13 \textperiodcentered{} Consecuencias, contrapartidas y mitigaciones}{tab:t130}%
  {F{0.26} F{0.26} Y}%
  {\cab{Consecuencia} & \cab{Trade-off} & \cab{Mitigación}}
  Dependencia de un servicio del proveedor de nube & Menor portabilidad del borde & Contratos en estándares abiertos; la lógica no vive en el gateway (criterio de reversibilidad) \\
  Costo por invocación en el peak de septiembre & Gasto variable & Cuotas por actor y límites declarados; modelado en FinOps con el peak de 105 TPS \\
  Los portales entran en enero de 2029 & La configuración crece en la Etapa 2 & Mismo gateway, rutas nuevas por configuración; sin componente adicional \\
\end{tablalafrox}

\section{ADR-14 --- Plataforma de Observabilidad Única: OTel + AMP + CloudWatch + X-Ray}
\label{sub:adr-14-plataforma-de-observabilidad-unica-}

Estado: Aprobado

Fecha: 2026-09-06

Decisión relacionada: D14 (Arquitectura Lógica (Subdocumento 4.1), Capa 8); reemplaza la parte de plataforma de D9.

Trazabilidad normativa: BA Art. 16.4 (\guillemotleft{}monitoreo del componente on-premise integrado a la misma plataforma de observabilidad que la nube, sin puntos ciegos\guillemotright{}) \textperiodcentered{} RT-03.16 (idéntica exigencia) \textperiodcentered{} RT-14.01 a RT-14.09 \textperiodcentered{} RT-09.01 (medición p95) \textperiodcentered{} Art. 16.3 (servicios administrados).

\paragraph{Contexto y problema.}
La observabilidad tiene que cubrir dos dominios muy distintos ---una nube elástica y cinco sitios que pueden quedar 24 h sin enlace--- y hacerlo sin puntos ciegos. La versión anterior de la arquitectura lógica declaraba un conjunto Prometheus + Grafana + Loki autoadministrado en cada centro de distribución, además de la plataforma en nube. Al contrastarlo con la vista física aparecieron dos problemas: ese conjunto no tenía máquina virtual dimensionada ---VM-06 es de 2 vCPU, 4 GB y 50 GB, insuficiente para sostener Prometheus, Loki y Grafana con 13 meses de métricas--- y, más de fondo, constituye una segunda plataforma de observabilidad, que es justamente lo que el Art. 16.4 y RT-03.16 prohíben al exigir \guillemotleft{}la misma\guillemotright{}.

\paragraph{Alternativas evaluadas.}
\begin{tablalafrox}{ADR-14 \textperiodcentered{} Alternativas evaluadas}{tab:t131}%
  {Y F{0.17} Y F{0.17} F{0.17}}%
  {\cab{Alternativa} & \cab{¿Una sola plataforma?} & \cab{Qué se ve durante un corte de 24 h} & \cab{Operación} & \cab{Veredicto}}
  A. Emisión on-premise (ADOT, buffer 24 h) $\rightarrow$ plataforma única en nube (AMP, CloudWatch, X-Ray, Grafana OSS) & Sí & Sin tableros centralizados; alarmas locales del equipamiento y bloqueo de frío 100 \% local; cero pérdida de telemetría gracias al buffer & Nula on-premise & Ganadora \\
  B. Prometheus + Grafana + Loki autoalojado por sitio, además de la nube & No: dos plataformas & Tableros locales completos & Dos plataformas que parchar, dimensionar y correlacionar, sobre un equipo de 4; requiere VM adicional no dimensionada & Rechazada (incumple RT-03.16 y Art. 16.4; sin emplazamiento) \\
  C. Datadog o New Relic & Sí & Depende del enlace igual que A & Menor, pero con lock-in y costo por host y por volumen & Rechazada (lock-in y costo) \\
\end{tablalafrox}

\paragraph{Criterio de selección y justificación técnico-económica.}
La norma pide una plataforma, no dos. Es el criterio decisivo: el Art. 16.4 y RT-03.16 usan la palabra \guillemotleft{}misma\guillemotright{}.

\guillemotleft{}Sin puntos ciegos\guillemotright{} se resuelve con el buffer, no con una segunda plataforma. El colector ADOT retiene 24 h en disco ---exactamente la autonomía comprometida del centro de distribución--- de modo que un corte no produce un hueco en la serie: produce un retraso que se cierra al reconectar.

Ninguna decisión crítica depende de un tablero. El bloqueo de despacho por excursión térmica es local (Decisión N\textdegree{} 4 del numeral 16.1 del caso), la alarma de cámara es acústica y luminosa, y los sensores de sala reportan al DCIM/BMS (RT-06.14). Lo que se pierde durante el corte es visibilidad agregada, no capacidad de operar.

Sin lock-in real: AMP es compatible con Prometheus y PromQL, de modo que las reglas de alerta y las consultas son portables; la instrumentación es OpenTelemetry, estándar neutral; y Grafana es OSS. Se conserva el ecosistema Prometheus/Grafana sin operar sus servidores.

Equipo de 4 personas: no se le entrega al CLIENTE una plataforma de observabilidad que mantener además del negocio.

\paragraph{Decisión adoptada.}
Instrumentación OpenTelemetry en todos los componentes; colectores ADOT on-premise (F-01: VM-06 Talca, VM-C04 Concepción, contenedor en cross-docking) con buffer en disco de 24 h; plataforma única en nube: métricas en AMP (13 meses), registros en CloudWatch Logs (12 meses en línea + 24 en archivo), trazas en X-Ray (30 días), tableros en Grafana OSS autoadministrado en sa-east-1 y alertas por AMP/Alertmanager y CloudWatch hacia SNS y PagerDuty. Se declara expresamente qué no está disponible durante un corte (RT-03.13), y se declara que ninguna decisión de la ventana 05:30--07:00 depende de ello.

\paragraph{Consecuencias, contrapartidas y mitigaciones.}
\begin{tablalafrox}{ADR-14 \textperiodcentered{} Consecuencias, contrapartidas y mitigaciones}{tab:t132}%
  {F{0.26} F{0.26} Y}%
  {\cab{Consecuencia} & \cab{Trade-off} & \cab{Mitigación}}
  Sin tableros centralizados durante un corte de enlace & Menor visibilidad agregada en contingencia & Buffer de 24 h sin pérdida; alarmas locales del equipamiento; bloqueo de frío local; declaración formal en la matriz RT-03.13 \\
  Dependencia de servicios del proveedor de nube & Portabilidad & AMP compatible con Prometheus/PromQL, OTel como instrumentación y Grafana OSS: el ecosistema es portable \\
  El buffer de 24 h consume disco en VM-06 y VM-C04 & Capacidad local & Dimensionado en el plan de capacidad; el corte de referencia es de 24 h (RT-03.10) \\
\end{tablalafrox}

\section{ADR-15 --- Gestión de Secretos: Servicio Administrado en lugar de Gestor Autoalojado}
\label{sub:adr-15-gestion-de-secretos-servicio-admini}

Estado: Aprobado

Fecha: 2026-09-06

Decisión relacionada: gestión de secretos de la Capa 7 de la arquitectura lógica (Subdocumento 4.1); se aplica en la sección de arquitectura de seguridad (subsección de cifrado y gestión de claves) y se emplaza como componente N-12 en la sección de modelo de emplazamiento híbrido.

Trazabilidad normativa: BA Art. 21.4 (\guillemotleft{}prohibición absoluta de credenciales, claves o secretos embebidos\dots{} uso obligatorio de un gestor de secretos con rotación automática\guillemotright{}) \textperiodcentered{} Art. 21.2 (gestión de claves con separación de funciones) \textperiodcentered{} Art. 16.2 (justificación de emplazamiento componente por componente) \textperiodcentered{} Art. 16.3 (servicios administrados) \textperiodcentered{} Art. 21 (Zero Trust) \textperiodcentered{} RT-04.09 \textperiodcentered{} RT-11.09.

\paragraph{Contexto y problema.}
La solución debe custodiar secretos de peso: credenciales del sistema de gestión de 2017, certificados AS2 y de intercambio electrónico de cada cadena de supermercados, credenciales de la autoridad tributaria y del procesador de pagos, y las credenciales de servicio entre módulos. El Art. 21.4 exige un gestor con rotación automática, y el Art. 16.2 exige que todo componente tenga emplazamiento justificado: un gestor de secretos no puede quedar fuera de la tabla de emplazamiento ni del dimensionamiento. La decisión es, por tanto, dónde vive el gestor y quién lo opera durante los 56 meses, con un equipo de TI del CLIENTE de cuatro personas.

\paragraph{Alternativas evaluadas.}
\begin{tablalafrox}{ADR-15 \textperiodcentered{} Alternativas evaluadas}{tab:t133}%
  {Y Y F{0.12} F{0.12} F{0.12} F{0.12}}%
  {\cab{Alternativa} & \cab{Emplazamiento y costo} & \cab{Operación para 4 personas} & \cab{Zero Trust} & \cab{Rotación automática} & \cab{Veredicto}}
  A. AWS Secrets Manager + SSM Parameter Store & Servicio administrado, sin VM ni licencia; costo por secreto declarable en FinOps & Nula: no se parcha ni se sella & Consumo saliente desde on-premise por VPC Endpoint; ningún puerto entrante & Nativa & Ganadora \\
  B. HashiCorp Vault autoalojado en Talca & VM adicional con alta disponibilidad, respaldo, procedimiento de sellado y desellado, y licencia empresarial si se requiere alta disponibilidad & Alta: el desellado tras un reinicio es un procedimiento manual crítico que puede caer a las 3 de la mañana & Equivalente & Sí, con configuración propia & Rechazada \\
  C. Secretos en variables de entorno o en la configuración del despliegue & --- & --- & --- & --- & Rechazada: prohibición absoluta del Art. 21.4 \\
\end{tablalafrox}

\paragraph{Criterio de selección y justificación técnico-económica.}
El componente debe existir en la vista física. Es la razón inmediata: lo declarado no estaba emplazado ni costeado. Cualquiera de las dos salidas corregía el defecto, pero solo una lo hacía sin agregar infraestructura.

Equipo de 4 personas (Cap. 2.4): el modo sellado de Vault tras un reinicio exige intervención humana con custodios. En una operación cuya ventana crítica es de 05:30 a 07:00 y cuyo turno de preparación es nocturno, introducir un componente que puede requerir desellado manual de madrugada es un riesgo operacional que no compra nada.

Zero Trust intacto: los nodos on-premise consumen secretos por VPC Endpoint saliente, coherente con la regla de que no se abren puertos entrantes salvo las dos excepciones D-AL-05.

La contingencia no depende de la nube. El secreto de la cuenta de emergencia queda deliberadamente fuera de línea, en sobre sellado con doble firma en el recinto de custodia de la sala (RT-06.26/06.27): es la única credencial que debe seguir siendo utilizable cuando no hay ni IdP ni enlace, y por eso no vive en ningún sistema.

Separación de funciones (Art. 21.2): quien administra la plataforma no administra las claves maestras; la política se expresa en IAM y queda auditada en CloudTrail.

\paragraph{Decisión adoptada.}
AWS Secrets Manager para secretos con rotación (credenciales del ERP, SII, Transbank y certificados AS2/EDI) y SSM Parameter Store para parámetros de configuración no sensibles, con consumo saliente desde los nodos on-premise por VPC Endpoint, rotación automática, y cifrado con CMK de KMS bajo separación de funciones. HashiCorp Vault queda declarado como alternativa evaluada y no adoptada. La cuenta de emergencia se custodia fuera de línea.

\paragraph{Consecuencias, contrapartidas y mitigaciones.}
\begin{tablalafrox}{ADR-15 \textperiodcentered{} Consecuencias, contrapartidas y mitigaciones}{tab:t134}%
  {F{0.26} F{0.26} Y}%
  {\cab{Consecuencia} & \cab{Trade-off} & \cab{Mitigación}}
  Los nodos on-premise requieren enlace para obtener un secreto nuevo & Dependencia de la nube para la rotación & Los secretos en uso se cachean en memoria por el proceso durante la autonomía de 24 h; ninguna rotación se programa dentro de la ventana crítica ni de un corte \\
  Dependencia de un servicio del proveedor de nube & Portabilidad & Los secretos son datos, no lógica; la exportación está cubierta por la declaración de reversibilidad (Art. 16.3) \\
  Costo por secreto y por rotación & Gasto variable & Volumen acotado (integraciones del catálogo, 15 entradas) y declarado en FinOps \\
\end{tablalafrox}

\section{Matriz de trazabilidad (ADR $\rightarrow$ normativa $\rightarrow$ decisión $\rightarrow$ documento)}
\label{sub:matriz-de-trazabilidad-adr-normativa-decis}

\begin{tablalafrox}{Trazabilidad de cada decisión de arquitectura con la normativa que la exige}{tab:t135}%
  {G{0.08} Y F{0.21} F{0.21} Y}%
  {\cab{ADR} & \cab{RT / RNF clave} & \cab{BA / Caso} & \cab{Decisión relacionada} & \cab{Documento que la materializa}}
  ADR-01 & RT-02.02 \textperiodcentered{} BTT sección 2.3 \textperiodcentered{} RT-09.05 \textperiodcentered{} RNF-19.01--04 & Art. 16 & D5 & Arquitectura de integración de este documento \textperiodcentered{} Dimensionamiento \\
  ADR-02 & RT-03.17 \textperiodcentered{} RT-03.10 \textperiodcentered{} RNF-13.01/13.07/13.08 \textperiodcentered{} RT-10.05 & Caso 02, numerales 6.12 y 8 & Decisión 16.1 N\textdegree{} 26 & Arquitectura de despliegue de este documento \textperiodcentered{} Emplazamiento \\
  ADR-03 & Art. 16.1--16.4 \textperiodcentered{} RT-03.01/03.02/03.10 \textperiodcentered{} RNF-02.01/05.01/05.02 & Art. 16 BA & D7 \textperiodcentered{} Decisión 16.1 N\textdegree{} 17 & Emplazamiento de este documento \textperiodcentered{} Despliegue \\
  ADR-04 & RT-05.02 \textperiodcentered{} RT-05.15 \textperiodcentered{} RNF-09.01/09.02/11.02 & Cap. 16.1 N\textdegree{} 18/35 \textperiodcentered{} D.S. 977/96 & D2/D4 & Persistencia de este documento \textperiodcentered{} Capa analítica y BI \\
  ADR-05 & RT-02.06/02.07 \textperiodcentered{} RT-03.10/03.12 \textperiodcentered{} RNF-07.01/13.01 & Art. 21 (Zero Trust) & D4 \textperiodcentered{} D-AL-03 & Arquitectura de integración de este documento \\
  ADR-06 & RT-12.10/12.11/12.12 \textperiodcentered{} RT-03.10 \textperiodcentered{} RNF-13.01 & Art. 22 (MFA) & D6 \textperiodcentered{} Decisión 16.1 N\textdegree{} 34 & Arquitectura de seguridad de este documento \\
  ADR-07 & RT-13.08 \textperiodcentered{} RT-12.11 \textperiodcentered{} RT-03.19 \textperiodcentered{} RNF-05.02/05.03/06.01 & Cap. 6/13.8 (caso) & Decisión de proyecto N\textdegree{} 19 & Emplazamiento de este documento \textperiodcentered{} RNF-05.02/05.03/06.01 \\
  ADR-08 & RT-05.11--05.15 \textperiodcentered{} RT-03.10 \textperiodcentered{} RNF-02.01 & Cap. 16.1 N\textdegree{} 14 & Decisión 16.1 N\textdegree{} 14 & Arquitectura de integración de este documento \textperiodcentered{} Emplazamiento \\
  ADR-09 & RT-07.01/07.07 \textperiodcentered{} RNF-20.06/20.07 & Art. 20 BA & Decisión 16.1 N\textdegree{} 27 & Arquitectura de despliegue de este documento \\
  ADR-10 & RT-03.14 \textperiodcentered{} RNF-13.03/13.04/13.05 & Cap. 16.1 N\textdegree{} 28 & Decisión 16.1 N\textdegree{} 28 & Dimensionamiento de este documento \textperiodcentered{} Emplazamiento (A-02) \\
  ADR-11 & RT-05.23 \textperiodcentered{} RT-16.16/16.17 \textperiodcentered{} RNF-12.01 \textperiodcentered{} RF-12.03/12.06/12.13 & Cap. 16.2 GS1 & RF-12 \textperiodcentered{} RF-01.10 & Arquitectura de integración de este documento \textperiodcentered{} RF-12.03/12.06/12.13 \\
  ADR-12 & RT-09.05 \textperiodcentered{} RT-03.06/03.08/03.09 \textperiodcentered{} RNF-19.01--04 & Cap. 14.2/15 \textperiodcentered{} RT-10.05 & Decisión 16.1 N\textdegree{} 30 & Dimensionamiento de este documento \\
  ADR-13 & RT-02.01 \textperiodcentered{} RT-11.11 \textperiodcentered{} RT-05.16/05.18 \textperiodcentered{} RT-02.02 & Art. 21.2 \textperiodcentered{} Art. 16.3 & D8 & Arquitectura de integración de este documento \textperiodcentered{} Arquitectura de seguridad de este documento \\
  ADR-14 & RT-03.16 \textperiodcentered{} RT-14.01--14.09 \textperiodcentered{} RT-09.01 & Art. 16.4 \textperiodcentered{} Art. 16.3 & D14 (reemplaza D9) & Arquitectura de seguridad de este documento \textperiodcentered{} RF-16.01 \\
  ADR-15 & RT-04.09 \textperiodcentered{} RT-11.09 \textperiodcentered{} RT-03.15 & Art. 21.4 \textperiodcentered{} Art. 21.2 \textperiodcentered{} Art. 16.2 \textperiodcentered{} Art. 16.3 & D13 \textperiodcentered{} SEC-01 & Arquitectura de seguridad de este documento \\
\end{tablalafrox}

LAFROX

Propuesta de Solución --- Distribuidora Puelche S.A.

Licitación N\textdegree{} TFEP-01/2026 --- Caso 02 Logística

SUBDOCUMENTO 5

Modelo y Gestión de Datos

Informe 1
